\chapter{Sesquilinear forms and quadratic forms}

\section{Sesquilinear forms}

\subsection{Bilinear maps.}

In this no., A and B denote two rings, E a left A-module, F a right B-module, and G an (A, B)-bimodule, that is, a commutative group equipped with a left A-module structure and a right B-module structure such that one has \((ag)b = a(gb)\) for all \(a \in A\) , \(b \in B\) , \(g \in G\) .

DEFINITION 1. --- One says that a map \(\Phi\) of the product \(\mathbf{E} \times \mathbf{F}\) in \(\mathbf{G}\) is bilinear if it satisfies the following conditions:

\begin{enumerate}
\def\labelenumi{(\arabic{enumi})}
\tightlist
\item
  \(\Phi(x + x', y) = \Phi(x, y) + \Phi(x', y)\)
\end{enumerate}

for all \(x \in \mathbf{E}\) , \(x' \in \mathbf{E}\) , \(y \in \mathbf{F}\) ;

\begin{enumerate}
\def\labelenumi{(\arabic{enumi})}
\setcounter{enumi}{1}
\tightlist
\item
  \(\Phi(x, y + y') = \Phi(x, y) + \Phi(x, y')\)
\end{enumerate}

for all \(x \in \mathbf{E}, y \in \mathbf{F}, y' \in \mathbf{F}\) ;

\begin{enumerate}
\def\labelenumi{(\arabic{enumi})}
\setcounter{enumi}{2}
\item
  \(\Phi(ax, y) = a\Phi(x, y)\) for all \(a \in \mathbf{A}\) , \(x \in \mathbf{E}\) , \(y \in \mathbf{F}\) ;
\item
  \(\Phi(x, yb) = \Phi(x, y)b\) for all \(x \in \mathbf{E}, y \in \mathbf{F}, b \in \mathbf{B}\) .
\end{enumerate}

The tensor product \(E \otimes_{Z} F\) is canonically endowed with a structure of (A, B)-bimodule characterized by \(a(x \otimes y)b = ax \otimes yb\) (Chap. III, 2nd ed., App. II, no. 3), and the data of a bilinear map \(\Phi\) of \(E \times F\) into G is equivalent to that of a map \(\Psi\) of \(E \otimes_{Z} F\) into G that is a homomorphism for the structures of (A, B)-bimodules and that satisfies \(\Psi(x \otimes y) = \Phi(x, y)\) for all \(x \in E\) and \(y \in F\) .

The conditions imposed on \(\Phi\) by Definition 1 mean that the partial maps \(d_{\Phi}(y): x \to \Phi(x, y)\) and \(s_{\Phi}(x): y \to \Phi(x, y)\) are respectively an A-linear map from E to G and a B-linear map from F to G. Let us equip the commutative group \(\mathcal{L}_{\mathbf{A}}(\mathbf{E}, \mathbf{G})\) (resp. \(\mathcal{L}_{\mathbf{B}}(\mathbf{F}, \mathbf{G})\) ) with the structure of a right B-module (resp. left A-module) defined by \(ub(x) = u(x). b\) ( \(u \in \mathcal{L}_{\mathbf{A}}(\mathbf{E}, \mathbf{G}), x \in \mathbf{E}, b \in \mathbf{B}\) ) (respectively \(av(y) = a. v(y)\) ( \(a \in \mathbf{A}, v \in \mathcal{L}_{\mathbf{B}}(\mathbf{F}, \mathbf{G}), y \in \mathbf{F}\) )). Then conditions (1) through (4) are respectively equivalent to:

\[
s _ {\Phi} (x + x ^ {\prime}) = s _ {\Phi} (x) + s _ {\Phi} (x ^ {\prime})\tag{1'}
\]

\[
d _ {\Phi} (y + y ^ {\prime}) = d _ {\Phi} (y) + d _ {\Phi} (y ^ {\prime})\tag{2'}
\]

\[
s _ {\Phi} (a x) = a. s _ {\Phi} (x)\tag{3'}
\]

\[
d _ {\Phi} (y b) = d _ {\Phi} (y). b,\tag{4'}
\]

for all \(x, x'\) in E, \(y, y'\) in F, \(a \in A\) , \(b \in B\) ; in other words, the map \(d_{\Phi}\) of F into \(\mathcal{L}_{A}(E, G)\) is B-linear, and the map \(s_{\Phi}\) from E into \(\mathcal{L}_{B}(F, G)\) is A-linear. We have, by definition,

\begin{enumerate}
\def\labelenumi{(\arabic{enumi})}
\setcounter{enumi}{4}
\tightlist
\item
  \(\Phi(x, y) = d_{\Phi}(y)(x) = s_{\Phi}(x)(y)\) for all \(x \in \mathbf{E}, y \in \mathbf{F}\) .
\end{enumerate}

DEFINITION 2. --- Given a bilinear map \(\Phi\) from E × F to G, the map \(d_{\Phi}\) of F into \(\mathfrak{L}_{\mathrm{A}}(\mathrm{E}, \mathrm{G})\) (resp. the map \(s_{\Phi}\) from E into \(\mathfrak{L}_{\mathrm{B}}(\mathrm{F}, \mathrm{G})\) ) characterized by (5) is called the linear map associated on the right (resp. on the left) to \(\Phi\) .

Conversely, the data of a B-linear map \(d\) of F into \(\mathfrak{L}_{\mathbf{A}}(\mathbf{E},\mathbf{G})\) (resp. of an A-linear map \(s\) from E into \(\mathfrak{L}_{\mathbf{B}}(\mathbf{F},\mathbf{G})\) ) determines uniquely, by the formula

\[
\Phi (x, y) = d (y) (x) \quad (\text { resp. } \Phi (x, y) = s (x) (y))
\]

a bilinear map \(\Phi\) of \(\mathbf{E} \times \mathbf{F}\) into \(\mathbf{G}\) , of which \(d\) (resp. \(s\) ) is the linear map associated on the right (resp. on the left).

DEFINITION 3. --- A bilinear map \(\Phi\) from E × F to G is said to be degenerate on the right (resp. on the left) if there exists a nonzero element \(y_0\) of F (resp. \(x_0\) of E) such that \(\Phi(x, y_0) = 0\) for all \(x \in \mathbf{E}\) , (resp. \(\Phi(x_0, y) = 0\) for all \(y \in \mathbf{F}\) ). One says that \(\Phi\) is degenerate if it is degenerate on the right or if it is degenerate on the left.

For \(\Phi\) to be nondegenerate on the right (resp. on the left) it is necessary and sufficient that the linear map associated on the right (resp. on the left) to \(\Phi\) be injective; to say that \(\Phi\) is nondegenerate therefore means that the associated linear maps \(d_{\Phi}\) and \(s_{\Phi}\) are both injective.

Let \((e_i)_{i\in I}\) and \((f_k)_{k\in K}\) be two families of elements of E and F, and let \((a_i)_{i\in I}\) and \((b_k)_{k\in K}\) be two families of elements of A and B, all but a finite number of which are zero. It follows from equalities (1) to (4), by induction on the number of nonzero coefficients, that one has

\[
\Phi (\sum_ {i} a _ {i} e _ {i}, \sum_ {k} f _ {k} b _ {k}) = \sum_ {i, k} a _ {i} \Phi (e _ {i}, f _ {k}) b _ {k}.\tag{6}
\]

If \((e_{i})\) and \((f_{k})\) are systems of generators of the modules E and F, \(\Phi\) is therefore completely determined by the elements \(g_{ik} = \Phi(e_{i}, f_{k})\) . If \((e_{i})\) and \((f_{k})\) are bases of E and F and elements \(g_{ik}\) of G \((i \in I, k \in K)\) are given, then the formula

\[
\Phi (\sum_ {i} a _ {i} e _ {i}, \sum_ {k} f _ {k} b _ {k}) = \sum_ {i, k} a _ {i} g _ {i k} b _ {k}\tag{6'}
\]

defines a map from E × F to G, which is bilinear and satisfies \(\Phi(e_{i}, f_{k}) = g_{ik}\) . When \((e_{i})\) and \((f_{k})\) are finite bases, we say that \((\Phi(e_{i}, f_{k}))\) is the matrix of \(\Phi\) with respect to these bases.

The bilinear maps from E × F into G obviously form a subgroup of the additive group of maps from E × F into G. On the other hand, let a (resp. b) be an element of the center of A (resp. B); then the map \(a\Phi b\) from E × F to G defined by \((a\Phi b)(x, y) = a \cdot \Phi(x, y) \cdot b\) is bilinear. The set of bilinear maps from E × F to G is thus endowed with a bimodule structure over the centers of A and B.

Let E' (resp. F') be a left A-module (resp. a right B-module), u (resp. v) a homomorphism from E to E' (resp. from F to F'), and \(\Phi'\) a bilinear map from \(E' \times F'\) into G. The inverse image of \(\Phi'\) (relative to u and v) is the bilinear map \(\Phi\) from E × F to G defined by

\[
\Phi (x, y) = \Phi^ {\prime} (u (x), v (y)) \quad (x \in \mathrm{E}, y \in \mathrm{F}).
\]

One verifies easily that one has

\[
d _ {\Phi} (y) = d _ {\Phi^ {\prime}} (\nu (y)) \circ u \quad \text { and } \quad s _ {\Phi} (x) = s _ {\Phi^ {\prime}} (u (x)) \circ \nu
\]

for all \(x \in \mathbf{E}, y \in \mathbf{F}\) .

Let \(\Phi\) be a bilinear map from \(\mathbf{E} \times \mathbf{F}\) into \(\mathbf{G}\) , and \(h\) a homomorphism (for the structures of (A, B)-bimodules) of \(\mathbf{G}\) into another (A, B)-bimodule \(\mathbf{G}'\) . Then \(h \circ \Phi\) is a bilinear map from \(\mathbf{E} \times \mathbf{F}\) into \(\mathbf{G}'\) .

\subsection{Sesquilinear maps.}

In this section, unless expressly stated otherwise, we denote by A and B two rings, by E a left A-module and by F a left B-module; we denote by \(b \to b^{\mathrm{J}} (b \in \mathrm{B})\) an anti-automorphism of B, that is, a bijection from B onto itself satisfying \((b + c)^{\mathrm{J}} = b^{\mathrm{J}} + c^{\mathrm{J}}\) and \((bc)^{\mathrm{J}} = c^{\mathrm{J}}b^{\mathrm{J}}\) for all \(b, c\) in B; we write \(\mathbf{J}'\) instead of \(\mathbf{J}^{-1}\) . We denote by G a \((\mathbf{A}, \mathbf{B})\) -bimodule (no. 1).

DEFINITION 4. --- A map \(\Phi\) from E \(\times\) F into G is said to be right-sesquilinear for J if it satisfies conditions (1), (2), (3) (def. 1, no. 1) as well as

\begin{enumerate}
\def\labelenumi{(\arabic{enumi})}
\setcounter{enumi}{6}
\tightlist
\item
  \(\Phi(x, by) = \Phi(x, y). b^J\) for all \(x \in \mathbf{E}, y \in \mathbf{F}\) and \(b \in \mathbf{B}\) .
\end{enumerate}

If J is the identity (which requires that B be commutative), we recover the notion of a bilinear map.

Let \((e_i)_{i\in I}\) and \((f_k)_{k\in K}\) be two families of elements of E and F, and let \((a_i)_{i\in I}\) and \((b_k)_{k\in K}\) be elements of A and B that are zero except for a finite number of them. We then have

\[
\Phi (\sum_ {i} a _ {i} e _ {i}, \sum_ {k} b _ {k} f _ {k}) = \sum_ {i, k} a _ {i} \Phi (e _ {i}, f _ {k}) b _ {k} ^ {\mathrm{J}}.\tag{8}
\]

As in the case of a bilinear map, the elements \(\Phi(e_i, f_k)\) determine \(\Phi\) in a unique way when \((e_i)\) and \((f_k)\) are systems of generators, and can be chosen arbitrarily when \((e_i)\) and \((f_k)\) are bases of E and F; when \((e_i)\) and \((f_k)\) are finite bases, we say that \((\Phi(e_i, f_k))\) is the matrix of \(\Phi\) with respect to these bases.

As for bilinear maps, one defines on the set of right sesquilinear maps (for J) from E × F into G a bimodule structure over the centers of A and B. One defines the notion of inverse image of a sesquilinear map by the same formula as for a bilinear map. We shall moreover see that the study of sesquilinear maps can be reduced to that of bilinear maps.

DEFINITION 5. --- Let B be a ring, F a left (resp. right) B-module and J an antiautomorphism of B. One denotes by F \(^{j}\) the right (resp. left) B-module having the same underlying additive group as F and in which the external composition law is \((b, y) \to b^{J'}y\) (resp. \((b, y) \to yb^{J'}\) ) \((b \in B, y \in F, J' = J^{-1})\) .

With the notations of Definition 5, a linear map from \(\mathbf{F}^{\mathrm{j}}\) into a right (resp. left) B-module H is therefore identified with a Z-linear map \(u\) of F into H satisfying

\[
u (b y) = u (y) b ^ {J} \quad (\text { resp. } u (y b) = b ^ {J} u (y)) \quad (b \in B, y \in F).
\]

The map \(u\) from F into H is a semilinear map of F into H relative to J (chap. II, App. I, no. 1), if one regards J as an isomorphism of the opposite ring B \(^{0}\) of B onto B, and F as a B \(^{0}\) -module on the right (resp. on the left).

Similarly, a right-sesquilinear map \(\Phi\) (for J) from E × F into G, where F is a left B-module, identifies with a bilinear map from E × F \(^{J}\) into G; if the latter is degenerate on the right (resp. degenerate on the left, nondegenerate), one says that \(\Phi\) is degenerate on the right (resp. degenerate on the left, nondegenerate).

Remark. --- Let A and B be two rings, \(J_{1}\) an anti-automorphism of A, M a right A-module, N a right B-module, and G an (A, B)-bimodule. One says that a map \(\Phi\) from M \(\times\) N into G is left-sesquilinear for J \(_{1}\) if it is Z-bilinear and if it satisfies

\[
\Phi (x a, y b) = a ^ {\mathrm{J} _ {1}} \Phi (x, y) b \quad (x \in \mathrm{M}, y \in \mathrm{N}, a \in \mathrm{A}, b \in \mathrm{B}).\tag{9}
\]

Such a map is identified with a bilinear map from \(M^{J_{1}} \times N\) into G. We will often leave it to the reader to transpose to left sesquilinear maps the definitions and properties given for right sesquilinear maps; when we speak of a sesquilinear map (without specifying), it will be a right sesquilinear map.

\subsection{Orthogonality. Direct sums of bilinear or sesquilinear maps.}

In this section, A and B denote rings, E a left A-module, F a right (resp. left) B-module, G an (A, B)-bimodule, and Φ a bilinear (resp. sesquilinear for a given antiautomorphism J of B) map from E × F to G.

DEFINITION 6. --- Two elements \(x \in E\) and \(y \in F\) are said to be orthogonal with respect to \(\Phi\) if \(\Phi(x, y) = 0\) . Two subsets \(E' \subset E\) and \(F' \subset F\) are said to be orthogonal if, for any \(x \in E'\) and \(y \in F'\) , \(x\) and \(y\) are orthogonal. The set of elements of E (resp. F) orthogonal to a given submodule N of F (resp. M of E) is a submodule of E (resp. F), called the totally orthogonal (or simply orthogonal) submodule to N (resp. M), and denoted \(N^0\) (resp. \(M^0\) ).

Let H and H' be two submodules of E or of F. One has \(\mathrm{H}\subset(\mathrm{H}^{0})^{0}\) (denoted \(H^{00}\) ); if \(H\subset H'\) , one has \(H^{\prime0}\supset H^{0}\) . It follows that one has \(H^{0}\supset(H^{00})^{0}\) and \(H^{0}\subset(H^{0})^{00}\) ; by setting

\[
\mathrm{H} ^ {0 0 0} = (\mathrm{H} ^ {0 0}) ^ {0} = (\mathrm{H} ^ {0}) ^ {0 0} = ((\mathrm{H} ^ {0}) ^ {0}) ^ {0},
\]

we therefore have \(\mathbf{H}^0 = \mathbf{H}^{000}\) .

For the map \(\Phi\) to be degenerate (no. 1, def. 3) it is necessary and sufficient that at least one of the two submodules E \(^{0}\) , F \(^{0}\) be \(\neq \{0\}\) . It is clear that \(\Phi(x, y)\) does not change when one adds to \(x\) (resp. \(y\) )

an element of F \(^{0}\) (resp. E \(^{0}\) ), and Φ therefore defines, by passing to the quotient, a bilinear (or sesquilinear) map on (E/F \(^{0}\) ) × (F/E \(^{0}\) ); this is visibly nondegenerate; it is called the nondegenerate bilinear (or sesquilinear) map associated with Φ.

Let \((\mathrm{E}_{i})_{i\in\mathrm{I}}\) be a family of left A-modules, \((\mathrm{F}_{i})_{i\in\mathrm{I}}\) a family of right B-modules (resp. left B-modules), \(\Phi_{i}\) a bilinear map (resp. sesquilinear on the right for J) from \(E_{i}\times F_{i}\) into G. Let E (resp. F) denote the direct sum module of the \(E_{i}\) (resp. \(F_{i}\) ). We see immediately that the map \(\Phi\) of \(E\times F\) into G defined by

\[
\Phi ((x _ {i}), (y _ {i})) = \sum_ {i} \Phi_ {i} (x _ {i}, y _ {i}) \quad (x _ {i} \in \mathrm{E} _ {i}, y _ {i} \in \mathrm{F} _ {i})\tag{10}
\]

(a sum that makes sense since its terms are zero except for finitely many of them) is bilinear (resp. sesquilinear on the right for J). It is called the direct sum of the maps \(\Phi_{i}\) . It is clear that \(\mathbf{E}_{i}\) is orthogonal to \(\mathbf{F}_{j}\) with respect to \(\Phi\) for \(i \neq j\) . Conversely, let \(\Phi\) be a bilinear or sesquilinear map from \(\mathbf{E} \times \mathbf{F}\) into G, and suppose that E is the direct sum of submodules \((\mathbf{E}_{i})_{i \in I}\) and F the direct sum of submodules \((\mathbf{F}_{i})_{i \in I}\) such that \(\mathbf{E}_{i}\) is orthogonal to \(\mathbf{F}_{j}\) for \(i \neq j\) ; then \(\Phi\) is the direct sum of its restrictions to the products \(\mathbf{E}_{i} \times \mathbf{F}_{i} (i \in I)\) .

For \(\Phi\) to be nondegenerate, it is necessary and sufficient that each of the \(\Phi_{i}\) be so; under these conditions, the submodule orthogonal to \(\mathbf{E}_{i}\) is \(\sum_{j\neq i}\mathbf{F}_{j}\) .

\subsection{Change of base rings.}

In this section, we denote by A, B, A', B' four rings, by h and h' homomorphisms from A to A' and from B to B' respectively, by G an (A, B)-bimodule, by G' an (A', B')-bimodule, and by u a homomorphism from the underlying abelian group of G to the underlying abelian group of G', satisfying

\[
u (a g b) = h (a) u (g) h ^ {\prime} (b) \quad (a \in \mathrm{A}, g \in \mathrm{G}, b \in \mathrm{B}).\tag{11}
\]

Let E (resp. F) be a left A-module (resp. a right B-module). Recall (Chap. III, 2nd ed., App. II, no. 10) that, if one regards A′ (resp. B′) as a right A-module (resp. left B-module), the tensor product \(E' = A' \otimes_{A} E\) (resp. \(F' = F \otimes_{B} B'\) ) is endowed with a left A'-module (resp. right B'-module) structure defined by

\[
\begin{array}{c c} a _ {1} ^ {\prime} (a ^ {\prime} \otimes x) = (a _ {1} ^ {\prime} a ^ {\prime}) \otimes x & (a ^ {\prime}, a _ {1} ^ {\prime} \in A ^ {\prime}, x \in E) \\ (\text { resp. } (y \otimes b ^ {\prime}) b _ {1} ^ {\prime} = y \otimes (b ^ {\prime} b _ {1} ^ {\prime}) & (b ^ {\prime}, b _ {1} ^ {\prime} \in B ^ {\prime}, y \in F)). \end{array}\tag{12}
\]

PROPOSITION 1. --- Let E be a left A-module and F a right B-module; set \(\mathbf{E}^{\prime} = \mathbf{A}^{\prime}\otimes_{\mathbf{A}}\mathbf{E}\) and \(\mathbf{F}^{\prime} = \mathbf{F}\otimes_{\mathbf{B}}\mathbf{B}^{\prime}\) . For every bilinear map \(\Phi\) of \(\mathbf{E}\times \mathbf{F}\) into G, there exists one and only one bilinear map \(\Phi^{\prime}\) from \(\mathbf{E}^{\prime}\times \mathbf{F}^{\prime}\) into G' such that one has

\[
\Phi^ {\prime} (a ^ {\prime} \otimes x, y \otimes b ^ {\prime}) = a ^ {\prime}. u (\Phi (x, y)). b ^ {\prime}\tag{13}
\]

for all \(a' \in \mathrm{A}'\) , \(b' \in \mathrm{B}'\) , \(x \in \mathrm{E}\) , \(y \in \mathrm{F}\) .

The uniqueness of \(\Phi'\) follows from the fact that the elements \(a' \otimes x\) and \(y \otimes b'\) generate E' and F' respectively. To prove their existence, consider the map

\[
m: (a ^ {\prime}, x, y, b ^ {\prime}) \rightarrow a ^ {\prime}. u (\Phi (x, y)). b ^ {\prime}
\]

from A' × E × F × B' into G; it is evidently Z-multilinear, and it satisfies

\[
\begin{array}{r} m (a ^ {\prime}, a x, y, b ^ {\prime}) = m (a ^ {\prime} h (a), x, y, b ^ {\prime}) \\ m (a ^ {\prime}, x, y b, b ^ {\prime}) = m (a ^ {\prime}, x, y, h ^ {\prime} (b) b ^ {\prime}) \\ (a \in \mathrm{A}, b \in \mathrm{B}, a ^ {\prime} \in \mathrm{A} ^ {\prime}, b ^ {\prime} \in \mathrm{B} ^ {\prime}, x \in \mathrm{E}, y \in \mathrm{F}). \end{array}
\]

There therefore exists a Z-bilinear map \(\Phi'\) from E' × F' into G' satisfying (13) (Chap. III, 2nd ed., App. II, no. 1, prop. 2). This relation and the definition of the module structures of E' and F' by (12) show that \(\Phi'\) is bilinear, which completes the proof.

Given the hypotheses and notation of Proposition 1, let us now study the linear maps associated with \(\Phi\) and with \(\Phi'\) (no. 1, def. 2). To this end, we shall first define a canonical homomorphism from \(\mathfrak{L}_{A}(\mathrm{E}, \mathrm{G})\) into \(\mathfrak{L}_{A'}(\mathrm{E}', \mathrm{G}')\) . For every \(\nu \in \mathfrak{L}_{\Lambda}(\mathrm{E}, \mathrm{G})\) the map \((a', x) \to a'. u(\nu(x))\) of \(\mathrm{A}' \times \mathrm{E}\) into \(\mathrm{G}'\) is Z-bilinear, and, in view of (11), maps \((a'h(a), x)\) and \((a', ax)\) ( \(a \in \mathrm{A}\) ) onto the same element of \(\mathrm{G}'\) ; it therefore defines (Chap. III, 2nd ed.,

App. II, nos. 1 and 10) a map \(k(\nu)\) of \(\mathbf{E}' = \mathbf{A}' \otimes_{\mathbf{A}} \mathbf{E}\) into \(\mathbf{G}'\) such that \(k(\nu)(a' \otimes x) = a'. u(\nu(x))\) , and which, by (12), is A'-linear. Moreover, one deduces immediately from (12) that the map \(\nu \to k(\nu)\) of \(\mathfrak{L}_{\mathbf{A}}(\mathbf{E}, \mathbf{G})\) into \(\mathfrak{L}_{\mathbf{A}'}(\mathbf{E}', \mathbf{G}')\) satisfies \(k(\nu b) = k(\nu) h'(b)\) for all \(b \in B\) . Let \(i\) denote the canonical map \(y \to y \otimes 1\) of F into \(\mathbf{F}'\) . Then the diagram

\[
\begin{array}{c c c} \mathrm{F} & \xrightarrow {d _ {\Phi}} & \mathcal {L} _ {A} (\mathrm{E}, \mathrm{G}) \\ \downarrow i & & \downarrow k \\ \mathrm{F} ^ {\prime} & \xrightarrow {d _ {\Phi^ {\prime}}} & \mathcal {L} _ {A^{\prime}} (\mathrm{E} ^ {\prime}, \mathrm{G} ^ {\prime}) \end{array}\tag{14}
\]

(where \(d_{\Phi}\) and \(d_{\Phi'}\) denote the linear maps associated on the right to \(\Phi\) and \(\Phi'\) ) is commutative. Indeed, for \(x \in \mathrm{E}, y \in \mathrm{F}\) and \(a' \in \mathrm{A}'\) , one has \(d_{\Phi'}(i(y))(a' \otimes x) = \Phi'(a' \otimes x, y \otimes 1) = a'. u(\Phi(x, y)) = a'. u(d_{\Phi}(y)(x))\) , that is to say \(d_{\Phi'}(i(y))(a' \otimes x) = k(d_{\Phi}(y))(a' \otimes x)\) . One has an analogous commutation relation for the linear maps \(s_{\Phi}\) and \(s_{\Phi'}\) associated on the left to \(\Phi\) and \(\Phi'\) .

PROPOSITION 2. --- Suppose that B and B' are equipped with anti-automorphisms J and I such that

\[
h ^ {\prime} (b ^ {J}) = h ^ {\prime} (b) ^ {I} \quad \text { for all } b \in \mathrm{B}.\tag{15}
\]

Let E be a left A-module and F a left B-module; set \(E' = A' \otimes_{A} E\) and \(F' = B' \otimes_{B} F\) . For every sesquilinear map (with respect to J) \(\Phi\) of \(E \times F\) into G, there exists one and only one sesquilinear map (for I) \(\Phi'\) from \(E' \times F'\) into \(G'\) such that

\[
\Phi^ {\prime} (a ^ {\prime} \otimes x, b ^ {\prime} \otimes y) = a ^ {\prime}. u (\Phi (x, y)). b ^ {\prime I}\tag{16}
\]

for all \(a' \in \mathbf{A}'\) , \(b' \in \mathbf{B}'\) , \(x \in \mathbf{E}\) , \(y \in \mathbf{F}\) .

The uniqueness of \(\Phi'\) follows from the fact that the tensor products \(a' \otimes x\) and \(b' \otimes y\) generate E′ and F′ respectively. To establish their existence, consider the map

\[
m: (a, x, b ^ {\prime}, y) \rightarrow a ^ {\prime}. u (\Phi (x, y)). b ^ {\prime I}
\]

of \(\mathbf{A}' \times \mathbf{E} \times \mathbf{B}' \times \mathbf{F}\) into \(\mathbf{G}\). It is evidently \(\mathbf{Z}\)-multilinear, and, taking (11) and (15) into account, satisfies \(m(a', ax, b', y) = m(a'h(a), x, b', y)\) and \(m(a', x, b', by) = a'. u(\Phi(x, y)) \cdot h'(b)^{I} b'^{I} = m(a', x, b'h'(b), y)\) ( \(a \in \mathbf{A}, b \in \mathbf{B}, a' \in \mathbf{A}'\) , \(b' \in \mathbf{B}'\) , \(x \in \mathbf{E}\) , \(y \in \mathbf{F}\) ). There therefore exists a \(\mathbf{Z}\)-bilinear map \(\Phi'\) from \(\mathbf{E}' \times \mathbf{F}'\) into \(\mathbf{G}'\)

satisfying (16) (chap. III, 2nd ed., App. II, no. 1, prop. 2). This relation, together with the definition of the module structures of E′ and F′ by (12), shows, in view of (15), that Φ′ is sesquilinear for I, which completes the proof.

The most important examples of \((\mathbf{A}^{\prime}, \mathbf{B}^{\prime})\) -bimodules \(G^{\prime}\) equipped with Z-linear maps u from G into \(G^{\prime}\) satisfying (11), are the following:

\begin{enumerate}
\def\labelenumi{\arabic{enumi})}
\tightlist
\item
  We take for G′ the tensor product A′ ⊗\_A G ⊗\_B B′ (chap. III, 2nd ed., App. II, no. 9) and for u the map
\end{enumerate}

\[
g \to 1 \otimes g \otimes 1\tag{\(g\in G\)}
\]

from G to G'. The pair (G', u) thus defined is visibly universal in the following sense: for every (A', B')-bimodule G₁' and every Z-linear map u₁ from G to G₁' satisfying the analogue of (11), there exists a unique Z-linear map f from G' to G₁' such that f(a'g'b') = a'f(g')b' (a' ∈ A', g' ∈ G', b' ∈ B'; in other words, f is a homomorphism for the bimodule structures of G' and G₁') and such that u₁ = f ∘ u.

\begin{enumerate}
\def\labelenumi{\arabic{enumi})}
\setcounter{enumi}{1}
\item
  When A = B = G (the structure of (A, A)-bimodule of A being defined by left and right homotheties), \(A' = B'\) , and \(h = h'\) one can take for \(G'\) the ring \(A'\) and for u the homomorphism h from A into \(A'\) .
\item
  Suppose that we have A = B, \(A' = B'\) , \(h = h'\) , that the rings A and \(A'\) are commutative, and that the structure of left A-module of G coincides with its structure of right A-module. One can then take for \(G'\) the tensor product \(A' \otimes_{A} G\) (the structure of \(A'\) -module on the right of \(G'\) coinciding with its structure of \(A'\) -module on the left) and for u the map \(g \to 1 \otimes g\) ( \(g \in G\) ) from G into \(G'\) . We shall then say that the bilinear (resp. sesquilinear) map \(\Phi'\) defined by prop. 1 (resp. prop. 2) is obtained from \(\Phi\) by extension of the base ring, or by extension of scalars.
\end{enumerate}

What follows is valid equally for bilinear maps and for sesquilinear maps; the hypotheses and notations are those of prop. 1 (resp. prop. 2). Given a submodule M of E or F, we denote by M' the submodule of E' or F' generated by the canonical image of M.

PROPOSITION 3. --- With the hypotheses and notations of prop. 1 (resp. prop. 2), suppose moreover that A, B, A', B' are fields and that the maps \(\alpha\) and \(\beta\) from A'⊗A G and G⊗B B' into G' characterized by \(\alpha(a'\otimes g)=a'u(g)\) and \(\beta(g\otimes b')=u(g)b'\) ( \(a'\in A'\) , \(b'\in B'\) , \(g\in G\) ) are injective. Let M be a subspace of E and N a subspace of F. Then the subspace \((M')^{0}\) of F' orthogonal to M' with respect to \(\Phi'\) is equal to \((M^{0})^{\prime}\) , and, similarly, we have \((N')^{0} = (N^{0})^{\prime}\) .

Indeed, the inclusions \((\mathbf{M}^{0})^{\prime}\subset(\mathbf{M}^{\prime})^{0}\) and \((\mathbf{N}^{0})^{\prime}\subset(\mathbf{N}^{\prime})^{0}\) are evident (and moreover true without hypotheses on A, B, A', B', \(\alpha\) or \(\beta\) ). We shall prove the inclusion \((\mathbf{M}^{\prime})^{0}\subset(\mathbf{M}^{0})^{\prime}\) ; we leave to the reader the verification of the inclusion \((\mathbf{N}^{\prime})^{0}\subset(\mathbf{N}^{0})^{\prime}\) , which is entirely analogous. Let \(y^{\prime}\) be an element of \((\mathbf{M}^{\prime})^{0}\) . One can write

\[
y ^ {\prime} = \sum_ {i = 1} ^ {s} y _ {i} \otimes b _ {i} ^ {\prime} \quad (\text { resp. } y ^ {\prime} = \sum_ {i = 1} ^ {s} b _ {i} ^ {\prime} \otimes y _ {i})
\]

where \(y_i \in \mathbf{F} (1 \leqslant i \leqslant s)\) , and where the \(b_i'\) are elements of \(\mathbf{B}'\) which are linearly independent over \(\mathbf{B}\) for the structure of \(\mathbf{B}\) -module on the left (resp. on the right) of \(\mathbf{B}'\) . Let \(x \in \mathbf{M}\) and \(x' = 1 \otimes x \in \mathbf{M}'\) . We have

\[
0 = \Phi^ {\prime} (x ^ {\prime}, y ^ {\prime}) = \sum_ {i} u (\Phi (x, y _ {i})) b _ {i} ^ {\prime} = \beta (\sum_ {i} \Phi (x, y _ {i}) \otimes b _ {i} ^ {\prime})
\]

\[
(\text { resp. } 0 = \Phi^ {\prime} (x ^ {\prime}, y ^ {\prime}) = \sum_ {i} u (\Phi (x, y _ {i})) b _ {i} ^ {\prime \mathrm{I}} = \beta (\sum_ {i} \Phi (x, y _ {i}) \otimes b _ {i} ^ {\prime \mathrm{I}}).
\]

Since β is injective and since the \(b_i'\) (resp. the \(b _ {i} ^ {\prime \mathrm{I}}\) , taking (15) into account) are linearly independent over B for the structure of left B-module of B', this entails \(\Phi(x, y_i) = 0\) for \(i = 1, \ldots, s\) . Since this last relation holds for every \(x \in \mathbf{M}\) , one has \(y_i \in \mathbf{M}^0\) for \(i = 1, \ldots, s\) , whence \(y' \in (\mathbf{M}^0)'\) . QED.

COROLLARY. --- With the hypotheses and notations of prop. 3, for \(\Phi'\) to be nondegenerate, it is necessary and sufficient that \(\Phi\) be nondegenerate.

Indeed, by prop. 3, we have \((\mathbf{F}')^0 = (\mathbf{F}^0)'\) and \((\mathbf{E}')^0 = (\mathbf{E}^0)'\) . On the other hand, for \(\Phi\) (resp. \(\Phi'\) ) to be nondegenerate, it is necessary and sufficient that one have \(\mathbf{F}^0 = \mathbf{E}^0 = \{0\}\) (resp. \((\mathbf{F}')^0 = (\mathbf{E}')^0 = \{0\}\) ).

Remark. --- Suppose that A, B, A' and B' are fields. Then, for the bimodules G' defined in the three examples above, the maps α and β are injective, as follows immediately from chap. III, 2nd ed., App. II, no. 6.

\subsection{Some identities.}

In this section, we denote by A a ring equipped with an anti-automorphism J, by E a left A-module, by G an (A, A)-bimodule, and by Φ a sesquilinear map (on the right) for J from E × E into G. We set Q(x) = Φ(x, x) (x ∈ E). We evidently have

\[
\mathrm{Q} (x + y) = \mathrm{Q} (x) + \Phi (x, y) + \Phi (y, x) + \mathrm{Q} (y)\tag{17}
\]

\[
\mathrm{Q} (x - y) = \mathrm{Q} (x) - \Phi (x, y) - \Phi (y, x) + \mathrm{Q} (y)\tag{18}
\]

whatever x, y in E. Hence, by subtraction,

\[
2 (\Phi (x, y) + \Phi (y, x)) = Q (x + y) - Q (x - y).\tag{19}
\]

Let \(a\) be an element of A; replacing \(y\) by \(ay\) in (19), we obtain

\[
2 (\Phi (x, y) a ^ {j} + a \Phi (y, x)) = Q (x + a y) - Q (x - a y).\tag{20}
\]

From (19) and (20), multiplying (19) by \(a\) (on the left) and subtracting, we obtain:

\[
\begin{array}{r l} 2 (a \Phi (x, y) - \Phi (x, y) a ^ {J}) & \\ = a Q (x + y) - a Q (x - y) - Q (x + a y) + Q (x - a y). \end{array}\tag{21}
\]

Suppose in particular that A is a quadratic extension K(i) of a commutative ring K, with \(i^{2} = -1\) (chap. II, § 7, no. 7), that J is the K-automorphism \(u + iv \rightarrow u - iv\) (u, v in K) of A, and that the left A-module and right A-module structures of G coincide. Setting a = i in (21), we obtain

\[
4 \Phi (x, y) = Q (x + y) - Q (x - y) + i Q (x + i y) - i Q (x - i y). \tag {22}
\]

\subsection{Bilinear and sesquilinear forms. Rank.}

In this section, A denotes a ring (resp. a ring equipped with an antiautomorphism J), E a left A-module, and F a right A-module (resp. a left A-module). One endows A with the structure of an (A, A)-bimodule defined by left multiplications and right multiplications. In this case, a bilinear (resp. right sesquilinear for J) map from E × F into the bimodule A is called a bilinear (resp. right sesquilinear for J) form on E × F.

When E = F (which implies that it is a sesquilinear form), one often says that a sesquilinear form on E × F is a sesquilinear form on E.

Given two left A-modules E and E′, and two sesquilinear forms \(\Phi\) and \(\Phi'\) for J on E and E' respectively, one says that \(\Phi\) and \(\Phi'\) are equivalent if there exists an isomorphism \(u\) of the A-module E onto the A-module E′ such that \(\Phi'(u(x), u(y)) = \Phi(x, y)\) for all \(x\) , \(y\) in E; then \(\Phi\) is the inverse image of \(\Phi'\) with respect to \(u\) and \(u\) , and \(\Phi'\) is the inverse image of \(\Phi\) with respect to \(u^{-1}\) and \(u^{-1}\) (no. 2).

Let \(\Phi\) be a bilinear form on \(E \times F\) (where F denotes a right A-module). The linear maps \(s_{\Phi}\) and \(d_{\Phi}\) associated with \(\Phi\) (no. 1, def. 2) are then maps from E into the dual \(F^{*}\) of F, and from F into the dual \(E^{*}\) of E.

Thus, by definition,

\[
\Phi (x, y) = \left\langle x, d _ {\Phi} (y) \right\rangle = \left\langle y, s _ {\Phi} (x) \right\rangle .\tag{23}
\]

We shall now define the linear maps associated with a sesquilinear form. Let J be an antiautomorphism of A, and \(\Phi\) a sesquilinear form (on the right) for J on E × F (where F denotes a left A-module); set J' = J \(^{-1}\) . The map \(\Phi'\) from F × E into A defined by

\[
\Phi^ {\prime} (y, x) = \Phi (x, y) ^ {J ^ {\prime}}
\]

\[
(x \in \mathrm{E}, y \in \mathrm{F})
\]

is, as is easily seen, a sesquilinear form (on the right) for J' on F × E. By no. 2 (def. 5), the sesquilinear forms Φ and Φ' identify respectively with bilinear forms on E × F' and on F × E'. The maps dΦ and dΦ' associated with the latter are called the maps associated on the right and on the left with the sesquilinear form \(\Phi\) , and are denoted \(d_{\Phi}\) and \(s_{\Phi}\) . We therefore have, by definition:

\[
\Phi (x, y) = \left\langle x, d _ {\Phi} (y) \right\rangle = \left\langle y, s _ {\Phi} (x) \right\rangle^ {J} \quad (x \in \mathrm{E}, y \in \mathrm{F}).\tag{24}
\]

Thus \(d_{\Phi}\) (resp. \(s_{\Phi}\) ) is a linear map from \(\mathbf{F}^{\mathrm{j}}\) into \(\mathbf{E}^{*}\) (resp. of \(\mathbf{E}^{\mathrm{j}'}\) into \(\mathbf{F}^{*}\) ), or equivalently a semilinear map from \(\mathbf{F}\) into \(\mathbf{E}^{*}\) (resp. of \(\mathbf{E}\) into \(\mathbf{F}^{*}\) ) relative to \(\mathbf{J}\) (resp. \(\mathbf{J}'\) ) if one regards \(\mathbf{J}\) (resp. \(\mathbf{J}'\) ) as a ring isomorphism \(\mathbf{A}^{0}\) (the opposite of \(\mathbf{A}\) ) on \(\mathbf{A}\) , and \(\mathbf{F}\) (resp. \(\mathbf{E}\) ) as a \(\mathbf{A}^{0}\) -module on the right.

Formula (24) and def. 6 of no. 3 immediately imply that for every submodule N of F (resp. M of E), we have

\[
\mathbf {N ^ {0}} = \stackrel {- 1} {s _ {\Phi}} (\mathbf {N ^ {\prime}}) \quad (\text { resp. } \mathbf {M ^ {0}} = \stackrel {- 1} {d _ {\Phi}} (\mathbf {M ^ {\prime}}))\tag{25}
\]

where N′ (resp. M′) is the submodule of the dual F* of F (resp. of the dual E* of E) orthogonal to N (resp. M) (Chap. II, § 4, no. 2).

PROPOSITION 4. --- Suppose that A is a field, and let \(\Phi\) be a bilinear form (resp. sesquilinear form with respect to J) on \(E \times F\) ; for \(E/F^{0}\) to be of finite dimension, it is necessary and sufficient that \(F/E^{0}\) be finite-dimensional, and these dimensions are then equal.

Indeed, let \(\Phi_{1}\) be the nondegenerate form associated with \(\Phi\) , on \((\mathrm{E} / \mathrm{F}^{0}) \times (\mathrm{F} / \mathrm{E}^{0})\) (no. 3). Suppose that \(\mathrm{E} / \mathrm{F}^{0}\) is of finite dimension \(n\) ; since the linear map \(d_{\Phi_1}\) of \(\mathrm{F} / \mathrm{E}^{0}\) (resp. \((\mathrm{F} / \mathrm{E}^{0})^{\mathrm{j}})\) into \((\mathrm{E} / \mathrm{F}^{0})^{*}\) is injective, \(\mathrm{F} / \mathrm{E}^{0}\) is of finite dimension \(n' \leqslant n\) ; by considering \(s_{\Phi_1}\) one sees likewise that \(n \leqslant n'\) .

COROLLARY 1. --- Suppose that A is a field and that \(\Phi\) is nondegenerate. For a subspace M of E to be of finite dimension, it is necessary and sufficient that M \(^{0}\) be of finite codimension in F, and one then has codim M \(^{0}\) = dim M, and M \(^{00}\) = M.

As \(F^{0}=\{0\}\) , the first two assertions follow from Prop. 4 applied to the restriction of \(\Phi\) to M × F. Moreover, \(M^{0}\) is the orthogonal of \(M^{00}\) , hence \(M^{00}\) is of finite dimension equal to codim \(M^{0}=dim M\) ; but since \(M^{00}\supset M\) , one has \(M^{00}=M\) .

COROLLARY 2. --- Under the hypotheses of cor. 1, let M, N be two subspaces of E; we then have \((M + N)^{0} = M^{0} \cap N^{0}\) ; if moreover M and N are of finite dimension, we have \((M \cap N)^{0} = M^{0} + N^{0}\) .

The first assertion is trivial. Suppose M and N are finite-dimensional, and let \(G = M^{0} + N^{0}\) ; we have \(G^{0} = M^{00} \cap N^{00} = M \cap N\) by cor. 1; prop. 4 applied to the restriction of \(\Phi\) to \(M \times G\) then shows (since \(M^{0} \subset G\) and \(G^{0} \subset M\) ) that one has dim \(M/(M \cap N) = \dim G/M^{0} = \text{codim } M^{0} - \text{codim } G\) , and since codim \(M^{0} = \dim M\) , we deduce dim \((M \cap N) = \text{codim } G\) . But we also have dim \((M \cap N) = \text{codim } (M \cap N)^{0}\) by cor. 1, and since \(G \subset G^{00} = (M \cap N)^{0}\) we have \(G = (M \cap N)^{0}\) .

Prop. 4 allows us to make the following definition:

DEFINITION 7. --- Let A be a field (resp. a field equipped with an antiautomorphism J), E a left vector space over A, F a right (resp. left) vector space over A, and Φ a bilinear (resp. sesquilinear for J) form on E × F. Suppose that E/F⁰ and F/E⁰ are of finite dimension over A. The rank of Φ is the (finite) common dimension of the vector spaces E/F⁰ and F/E⁰.

When \(E/F^{0}\) and \(F/E^{0}\) are of infinite dimension, one says that \(\Phi\) is of infinite rank.

PROPOSITION 5. --- With the hypotheses and notation as in def. 7, the linear maps \(s_{\Phi}\) and \(d_{\Phi}\) associated with \(\Phi\) have the same rank, and this rank is equal to the rank of the form \(\Phi\) .

Indeed, the kernel of the map \(d_{\Phi}\) of F into E* is evidently E°, hence its rank equals the dimension of F/E°. Likewise, the rank of \(s_{\Phi}\) is equal to the dimension of E/F°.

PROPOSITION 6. --- With the hypotheses and notation of Definition 7, suppose further that E and F have the same finite dimension. Then the following conditions are equivalent:

\begin{enumerate}
\def\labelenumi{\alph{enumi})}
\item
  \(d_{\Phi}\) is injective;
\item
  \(d_{\Phi}\) is surjective;
\item
  \(s_{\Phi}\) is injective;
\end{enumerate}

\(d)\) \(s_{\Phi}\) is surjective;

\begin{enumerate}
\def\labelenumi{\alph{enumi})}
\setcounter{enumi}{4}
\tightlist
\item
  \(\Phi\) is nondegenerate.
\end{enumerate}

Indeed, since E, F, E* and F* have the same finite dimension, a) and b) are equivalent, as are c) and d) (chap. II, § 3, no. 4). Since \(s_{\Phi}\) and \(d_{\Phi}\) have the same rank (prop. 5), a) and c) are equivalent. Since e) is equivalent to the relation \(\mathrm{E}^{0} = \mathrm{F}^{0} = \{0\}\) , it is equivalent to the conjunction of a) and c), whence the equivalence of the stated conditions.

COROLLARY. --- With the hypotheses and notations being those of def. 7, suppose further that E is finite-dimensional and that \(\Phi\) is nondegenerate. Then we have dim E = dim F and, for every basis \((e_i)\) ( \(1 \leqslant i \leqslant \dim E\) ) of E, there exists a basis \((f_i)\) of F such that \(\Phi(e_i, f_k) = \delta_{ik}\) ( \(i, k = 1, \ldots, \dim E\) ).

Indeed, as \(\Phi\) is nondegenerate, we have \(E^{0} = F^{0} = \{0\}\) , whence dim \(E = \dim F\) (prop. 4). It follows (prop. 6) that \(d_{\Phi}\) is an isomorphism of \(F\) (resp. \(F^{j}\) ) onto \(E^{*}\) ; therefore, if \((e_{i}^{*})\) is the dual basis of \((e_{i})\) , the elements \(f_{i} = d_{\Phi}^{-1}(e_{i}^{*})\) form a basis of \(F\) which, in view of formula (23) (resp. formula (24)), satisfies \(\Phi(e_{i}, f_{k}) = \delta_{ik}\) .

It is immediate that, in this corollary, one may interchange the roles of E and F, replacing \(d_{\Phi}\) by \(s_{\Phi}\) in the proof.

Remark. --- Let A be a ring equipped with an antiautomorphism J, and let M and N be right A-modules, and \(\Phi\) a left sesquilinear form for J on M × N (no. 2, Remark); it therefore satisfies the equality

\[
\Phi (x a, y a ^ {\prime}) = a ^ {j} \Phi (x, y) a ^ {\prime} \quad (a, a ^ {\prime} \in A, x \in M, y \in N).
\]

The map \(\Phi'\) from N × M into A defined by \(\Phi'(y, x) = \Phi(x, y)^{J'}\) (where J' = J \(^{-1}\) ) is a left sesquilinear form for J', and \(\Phi\) and \(\Phi'\) are identified with bilinear forms on M \(^{J'}\) × N and N \(^{J'}\) × M, respectively. The maps \(s_{\Phi}\) and \(s_{\Phi'}\) associated with these bilinear forms are called the maps associated on the left and on the right with the sesquilinear form \(\Phi\) , and are denoted \(s_{\Phi}\) and \(d_{\Phi}\) . We therefore have, by definition

\[
\Phi (x, y) = \left\langle y, s _ {\Phi} (x) \right\rangle = \left\langle x, d _ {\Phi} (y) \right\rangle^ {J} \quad (x \in M, y \in N),\tag{26}
\]

and \(s_{\Phi}\) (resp. \(d_{\Phi}\) ) is a linear map from \(\mathbf{M}^{\mathbf{j}}\) into \(\mathbf{N}^{*}\) (resp.

of N \(^{j'}\) into M*). One could easily state and prove the analogues, for the case considered here, of def. 7 and Props. 4, 5, 6.

\subsection{Inverse form of a bilinear or sesquilinear form.}

Let A be a ring, E a left A-module, F a right A-module, and \(\Phi\) a bilinear form on \(\mathbf{E} \times \mathbf{F}\) . It is assumed here that the associated maps of \(\Phi\) , which will be denoted \(s\) and \(d\) , are bijective. Then the product map \((s, d)\) is a bijection from \(\mathbf{E} \times \mathbf{F}\) onto \(\mathbf{F}^* \times \mathbf{E}^*\) , and defines, by transport of structure, a bilinear form \(\hat{\Phi}\) on \(\mathbf{F}^* \times \mathbf{E}^*\) . This therefore satisfies

\[
\begin{array}{r l} \hat {\Phi} (y ^ {\prime}, x ^ {\prime}) = & \Phi (s ^ {- 1} (y ^ {\prime}), d ^ {- 1} (x _ {i} ^ {\prime})) \\ = & \left\langle s ^ {- 1} (y ^ {\prime}), x ^ {\prime} \right\rangle = \left\langle d ^ {- 1} (x ^ {\prime}), y ^ {\prime} \right\rangle \qquad (x ^ {\prime} \in \mathrm{E} ^ {*}, y ^ {\prime} \in \mathrm{F} ^ {*}). \end{array}\tag{27}
\]

DEFINITION 8. --- Let \(\Phi\) be a bilinear form on E × F whose associated maps s and d are bijective. The bilinear form \(\hat{\Phi}\) on F* × E* defined by (27) is called the inverse form of \(\Phi\) .

Now let \(\hat{s}\) and \(\hat{d}\) be the linear maps from F* to E** and from E* to F** associated on the left and on the right to \(\hat{\Phi}\) . Since, for \(x' \in \mathrm{E}^*\) and \(y' \in \mathrm{F}^*\) , we have by definition

\[
\hat {\Phi} (y ^ {\prime}, x ^ {\prime}) = \left\langle y ^ {\prime}, \hat {d} (x ^ {\prime}) \right\rangle = \left\langle x ^ {\prime}, \hat {s} (y ^ {\prime}) \right\rangle
\]

Comparing with (27), one sees that the linear form \(\hat{d}(x')\) on F* is equal to that defined by the element \(d^{-1}(x')\) of F. It follows that the composite map \(\hat{d} \circ d\) is the canonical map from F to its bidual F**, and that this map is bijective since \(d\) and \(\hat{d}\) (the latter by transport of structure) are bijective; hence, if one canonically identifies F with F**, one has \(\hat{d} = d^{-1}\) . Similarly, E is canonically identified with E**, and the canonical map from E to E** is \(\hat{s} \circ s\) and one has \(\hat{s} = s^{-1}\) . It follows from this that the inverse form of \(\hat{\Phi}\) is \(\Phi\) .

Now consider a ring A equipped with an antiautomorphism J, two left A-modules E and F, and a right sesquilinear form \(\Phi\) for J on E \(\times\) F such that the associated maps of \(\Phi\) , which will be denoted \(s\) and \(d\) , are bijective. Define a map \(\hat{\Phi}\) of \(\mathbf{F}^{*} \times \mathbf{E}^{*}\) into A by the first equation (27). This map satisfies, by (24) (no. 6), the relation

\[
\hat {\Phi} (y ^ {\prime}, x ^ {\prime}) = \left\langle s ^ {- 1} (y ^ {\prime}), x ^ {\prime} \right\rangle = \left\langle d ^ {- 1} (x ^ {\prime}), y ^ {\prime} \right\rangle^ {J} \quad (x ^ {\prime} \in E ^ {*}, y ^ {\prime} \in F ^ {*}).\tag{28}
\]

The map \(\hat{\Phi}\) is evidently \(\mathbf{Z}\) -bilinear; moreover, we have, for \(a\) , \(b\) in A, \(x' \in \mathrm{E}^*\) and \(y' \in \mathrm{F}^*\) and by virtue of the definitions of \(s\) and \(d\) ,

\[
\widehat {\Phi} (y ^ {\prime} a, x ^ {\prime} b) = \Phi (a ^ {j} s ^ {- 1} (y ^ {\prime}), b ^ {j ^ {\prime}} d ^ {- 1} (x ^ {\prime})) = a ^ {j} \widehat {\Phi} (y ^ {\prime}, x ^ {\prime}) b;
\]

therefore \(\hat{\Phi}\) is a left sesquilinear form for J (no. 2) on \(\mathbf{F}^{*} \times \mathbf{E}^{*}\) .

DEFINITION 9. --- Let \(\Phi\) be a right-sesquilinear form for J on E × F, whose associated maps s and d are bijective. The left-sesquilinear form \(\hat{\Phi}\) for J on F* × E* is called the inverse form of \(\Phi\) .

We leave it to the reader to define and study the inverse form of a left sesquilinear form. This inverse form is a right sesquilinear form.

Let \(\hat{s}\) and \(\hat{d}\) be the associated maps of \(\hat{\Phi}\) ; by (26) (no. 6) we have

\[
\hat {\Phi} (y ^ {\prime}, x ^ {\prime}) = \left\langle y ^ {\prime}, \hat {d} (x ^ {\prime}) \right\rangle^ {\mathbf {J}} = \left\langle x ^ {\prime}, \hat {s} (y ^ {\prime}) \right\rangle .\tag{29}
\]

Because of the fact that \(s\) is bijective, and from the equality \(\langle s^{-1}(y'), x'\rangle = \langle y', \hat{d}(x') \rangle^{\mathfrak{J}}\) which results from (28) and (29), we deduce that \(\hat{d}\) is bijective; therefore \(\hat{d} \circ d\) is bijective. But the equality \(\langle d^{-1}(x'), y' \rangle = \langle y', \hat{d}(x') \rangle\) which results from (28) and (29), shows that \(\hat{d} \circ d\) is the canonical map from F into its bidual F**. One sees similarly that \(\hat{s}\) is bijective and that \(\hat{s} \circ s\) is the canonical map from E into E**. Therefore, if one identifies E** with E and F** with F by means of these canonical maps, we have \(\hat{s} = s^{-1}\) , \(\hat{d} = d^{-1}\) , and \(\Phi\) is the inverse form of \(\hat{\Phi}\) .

With the same notations and hypotheses, let a be an invertible element of the center of A. Then the maps associated with the form \(a\Phi\) are, by (23) (resp. (24)), equal to \(a.d\) and \(a.s\) (resp. \(a^{j'}\) .s), hence are bijective. It thus follows from (27) that the inverse form of \(a\Phi\) is \(a^{-1}\widehat{\Phi}\) (resp. \((a^{j'})^{-1}\widehat{\Phi}\) ).

\subsection{Adjoint of a homomorphism.}

In this section, we denote by A a ring (resp. a ring equipped with an antiautomorphism J), by E and E' two left A-modules, by F and F' two right A-modules (resp. left A-modules), and by Φ and Φ' two bilinear forms (resp. sesquilinear forms for J) on E × F and E' × F' respectively. We assume that Φ is nondegenerate, in other words (no. 1) that the linear maps \(d_{\Phi}\) and \(s_{\Phi}\) associated with Φ are injective.

Given a homomorphism \(u\) from E into \(\mathbf{E}'\) , consider the set \(\mathbf{F}_{1}'\) of elements \(y'\) of \(\mathbf{F}'\) such that there exists \(y \in \mathbf{F}\) for which one has \(d_{\Phi'}(y') \circ u = d_{\Phi}(y)\) , that is to say \(\Phi'(u(x), y') = \Phi(x, y)\) for all \(x \in \mathbf{E}\) . It is clear that \(\mathbf{F}_{1}'\) is a submodule of \(\mathbf{F}'\) . Since \(d_{\Phi}\) is injective, there exists, for every \(y' \in \mathbf{F}_{1}'\) , one and only one element \(y\) of F such that \(\Phi'(u(x), y') = \Phi(x, y)\) . The map \(y' \to y\) of \(\mathbf{F}_{1}'\) into F thus defined is A-linear; denoting it by \(u^{*}\) , one has, for every \(x \in \mathbf{E}\) and every \(y' \in \mathbf{F}_{1}'\)

\[
\Phi^ {\prime} (u (x), y ^ {\prime}) = \Phi (x, u ^ {*} (y ^ {\prime})).\tag{30}
\]

DEFINITION 10. --- The hypotheses and notations being as before, one says that the homomorphism \(u^*\) of \(F_1'\) into F satisfying (30) is the left adjoint of \(u\) , and that \(F_1'\) is the defining submodule of \(u^*\) .

One defines similarly the right adjoint of a homomorphism \(\nu\) from F into F' by the formula

\[
\Phi^ {\prime} (x ^ {\prime}, \nu (y)) = \Phi (\nu^ {*} (x ^ {\prime}), y) \quad (x ^ {\prime} \in \mathrm{E} _ {1} ^ {\prime}, y \in \mathrm{F}),\tag{31}
\]

where \(E_{1}^{\prime}\) denotes the submodule of \(E^{\prime}\) defined analogously to \(F_{1}^{\prime}\) .

Remark. --- If the left adjoint \(u^{*}\) of \(u: E \to E'\) is everywhere defined, and if \(s_{\Phi'}\) and \(d_{\Phi'}\) are injective, formula (30) shows that \(u\) is the right adjoint of \(u^{*}\) .

From (30), we deduce that, if \(u_1\) and \(u_2\) are two homomorphisms from E into E′ admitting everywhere-defined adjoints, and if c is an element of the center of A, then we have

\[
\left\{ \begin{array}{l} (u _ {1} + u _ {2}) ^ {*} = u _ {1} ^ {*} + u _ {2} ^ {*}; 1 ^ {*} = 1; \\ (c u _ {1}) ^ {*} = c. u _ {1} ^ {*} \quad \text { when } \Phi \text { and } \Phi^ {\prime} \text { are bilinear }; \\ (c u _ {1}) ^ {*} = c ^ {J ^ {\prime}}. u _ {1} ^ {*} \quad \text { when } \Phi \text { and } \Phi^ {\prime} \text { are sesquilinear }. \end{array} \right.\tag{32}
\]

Additionally, if \(E^{\prime\prime}\) is a third left A-module, \(F^{\prime\prime}\) a third right (resp. left) A-module, \(\Phi^{\prime\prime}\) a bilinear form (resp. sesquilinear form with respect to J) on \(E^{\prime\prime} \times F^{\prime\prime}\) , and if \(u^{\prime}\) is a homomorphism of \(E^{\prime}\) into \(E^{\prime\prime}\) admitting a (left) adjoint defined everywhere, we have

\[
(u ^ {\prime} \circ u) ^ {*} = u ^ {*} \circ u ^ {\prime *}.\tag{33}
\]

In particular, if \(u\) is an isomorphism from E onto E', and if the adjoints \(u^{*}\) and \((u^{-1})^{*}\) are everywhere defined, \(u^{*}\) is an isomorphism from F' onto F, and we have \((u^{*})^{-1} = (u^{-1})^{*}\) . Analogous properties hold for right adjoints.

PROPOSITION 7. --- With the same notation as above, assume that \(d_{\Phi}\) is bijective. Then every homomorphism \(u\) from E to E' admits an everywhere-defined left adjoint, and we have \(u^{*} = (d_{\Phi})^{-1} \circ {}^{t}u \circ d_{\Phi'}\) .

Indeed, as \(d_{\Phi}\) is bijective, we have, with the notations from the beginning of the paragraph, \(\mathbf{F}_{1}^{\prime}=\mathbf{F}^{\prime}\) , and \(u^{*}\) is therefore everywhere defined. On the other hand, (30) is equivalent to

\[
\left\langle u (x), d _ {\Phi^ {\prime}} \left(y ^ {\prime}\right) \right\rangle = \left\langle x, \left(d _ {\Phi} \circ u ^ {*}\right) \left(y ^ {\prime}\right) \right\rangle \quad \left(x \in \mathrm{E}, y ^ {\prime} \in \mathrm{F} ^ {\prime}\right);
\]

or \(\langle d_{\Phi'}(y'), u(x) \rangle = \langle {}^t u(d_{\Phi'}(y')), x \rangle\) ; one therefore has \({}^t u(d_{\Phi'}(y')) = d_{\Phi}(u^*(y'))\) for all \(y' \in F'\) , whence \({}^t u \circ d_{\Phi'} = d_{\Phi} \circ u^*\) , and consequently the announced expression of \(u^*\) . QED.

Remark. --- When \(s_{\Phi}\) is bijective, every homomorphism \(\nu\) from F to F' admits a right adjoint that is everywhere defined, and we have

\[
v ^ {*} = (s _ {\Phi}) ^ {- 1} \circ {} ^ {t} v \circ s _ {\Phi^ {\prime}}.\tag{34}
\]

PROPOSITION 8. --- With the same notation as above, assume that \(s_{\Phi}\) and \(d_{\Phi}\) are bijective. Let \(u\) and \(v\) be isomorphisms from E onto \(\mathbf{E}'\) and from F onto \(\mathbf{F}'\) respectively. Then, for \(\Phi\) to be the inverse image of \(\Phi'\) with respect to \(u\) and \(\nu\) (that is, so that one has \(\Phi(x,y)=\Phi'(u(x),\nu(y))\) for all \(x\in\mathbf{E},y\in\mathbf{F}\) ), it is necessary and sufficient that we have \(u^{-1}=\nu^{*}\) and \(\nu^{-1}=u^{*}\) .

Indeed \(\Phi'(u(x), v(y)) = \Phi(x, y)\) is also written \(\Phi(x, u^*(v(y))) = \Phi(x, y)\) . If this holds for all \(x \in E\) and \(y \in F\) , one has \(u^* \circ v = 1\) since \(\Phi\) is nondegenerate. We therefore also have \(v^* \circ u = 1\) by (33). The converse is immediate.

COROLLARY. --- Let A be a ring equipped with an antiautomorphism J, and let E be a left A-module, \(\Phi\) a sesquilinear form for J on E × E whose associated maps are bijective, and u an automorphism of the A-module E. For u to leave \(\Phi\) invariant (that is, so that we have \(\Phi(u(x), u(y)) = \Phi(x, y)\) for all x, y in E), it is necessary and sufficient that the two adjoints of u be equal and that we have \(u^{*} = u^{-1}\) .

This follows immediately from prop. 8.

Remark. --- Under the hypotheses of the corollary to Proposition 8, suppose further that A is a field and that E is finite-dimensional over A. Let \(w\) be an endomorphism of E, \(w_{1}\) and \(w_{2}\) its right and left adjoints. Each of the conditions \(ww_{1}=1\) , \(ww_{2}=1\) , \(w_{1}w=1\) , \(w_{2}w=1\) implies that \(w\) is an automorphism of E leaving \(\Phi\) invariant, and that \(w_{1}=w_{2}\) .

\subsection{Tensor products and exterior powers of sesquilinear forms.}

In this section, A denotes a commutative ring. A bilinear form on a product of two A-modules is therefore a special case of a sesquilinear form. We denote by J an automorphism of A, and by J' its inverse.

Let \(\mathbf{E}_i\) ( \(i = 1, \ldots, m\) ) be A-modules. The map

\[
(x _ {1}, \dots , x _ {m}) \rightarrow x _ {1} \otimes \dots \otimes x _ {m}
\]

of \(\prod_{i=1}^{m}\mathrm{E}_{i}^{j}\) into \((\bigotimes_{i=1}^{m}\mathrm{E}_{i})^{j}\) ( \(x_{i}\in\mathrm{E}_{i}^{j}\) ) (cf. def. 5, no. 2) is evidently A-multilinear; it therefore defines (chap. III, § 1, no. 7) an A-linear map \(f\) of \(\bigotimes_{i}\mathrm{E}_{i}^{\mathrm{J}}\) into \((\bigotimes_{i}\mathrm{E}_{i})^{\mathrm{J}}\) ; this map sends \(x_{1}\otimes\cdots\otimes x_{m}\) (where the symbols \(\otimes\) denote tensor products in \(\bigotimes_{i}\mathrm{E}_{i}^{\mathrm{J}}\) ) to \(x_{1}\otimes\cdots\otimes x_{m}\) (where the symbols \(\otimes\) denote tensor products in \((\bigotimes_{i}\mathrm{E}_{i})^{\mathrm{J}}\) ). Hence \(f\) is an isomorphism of \(\bigotimes_{i}\mathrm{E}_{i}^{\mathrm{J}}\) onto \((\bigotimes_{i}\mathrm{E}_{i})^{\mathrm{J}}\) . We shall identify these two modules by means of this isomorphism.

Similarly, let E be an A-module. The map

\[
(x _ {1}, \dots , x _ {m}) \rightarrow x _ {1} \wedge \dots \wedge x _ {m}
\]

of \((\mathrm{E}^{\mathfrak{J}})^{m}\) into \((\wedge^{\mathfrak{m}}\mathrm{E})^{\mathfrak{J}}\) is obviously A-multilinear and alternating. It therefore defines an A-linear map \(f\) of \(\wedge^{\mathfrak{m}}\mathrm{E}^{\mathfrak{J}}\) into \((\wedge^{\mathfrak{m}}\mathrm{E})^{\mathfrak{J}}\) , which is evidently an isomorphism. We shall identify \(\wedge^{\mathfrak{m}}\mathrm{E}^{\mathfrak{J}}\) and \((\wedge^{\mathfrak{m}}\mathrm{E})^{\mathfrak{J}}\) by means of this isomorphism.

Let \(x'\) be an element of the dual \(\mathbf{E}^*\) of \(\mathbf{E}\) . The map \(x \to \langle x, x' \rangle^{\mathfrak{J}}\) ( \(x \in \mathbf{E}\) ) is an element \(x'^{\mathfrak{J}}\) of \((\mathrm{E}^{\mathfrak{J}})^*\) , and it is immediate that \(x' \to x'^{\mathfrak{J}}\) is a bijection \(g\) of \(\mathbf{E}^*\) onto \((\mathrm{E}^{\mathfrak{J}})^*\) satisfying \(g(ax') = a^{\mathfrak{J}}g(x')\) for all \(a \in \mathbf{A}\) . Consequently the composite map of \(g\) and of the identity map of \((\mathrm{E}^*)^{\mathfrak{J}}\) onto \(\mathbf{E}^*\) is an isomorphism of \((\mathrm{E}^*)^{\mathfrak{J}}\) onto \((\mathrm{E}^{\mathfrak{J}})^*\) . We shall identify these modules by means of this isomorphism, and we shall denote them \(\mathbf{E}_j^*\) .

Let \(E_{i}, F_{i} (i = 1, \ldots, m)\) be A-modules, and \(\Phi_{i} (i = 1, \ldots, m)\) a sesquilinear form for J on \(E_{i} \times F_{i}\) . The map

\[
(x _ {1}, \dots , x _ {m}, y _ {1}, \dots , y _ {m}) \rightarrow \Phi_ {1} (x _ {1}, y _ {1}) \Phi_ {2} (x _ {2}, y _ {2}) \dots \Phi_ {m} (x _ {m}, y _ {m})
\]

\((x_{i} \in \mathrm{E}_{i}, y_{i} \in \mathrm{F}_{i}, i = 1, \ldots, m)\) is an A-multilinear map from \(\mathrm{E}_{1} \times \cdots \times \mathrm{E}_{m} \times \mathrm{F}_{1}^{J} \times \cdots \times \mathrm{F}_{m}^{J}\) into A, and thus defines a bilinear form on \((\bigotimes_{i} \mathrm{E}_{i}) \times (\bigotimes_{i} \mathrm{F}_{i}^{J})\) (chap. III, § 1, no. 7). Since the second factor has been identified with \((\bigotimes_{i} \mathrm{F}_{i})^{J}\) , we have therefore defined a sesquilinear form \(\Phi\) for J on \((\bigotimes_{i} \mathrm{E}_{i}) \times (\bigotimes_{i} \mathrm{F}_{i})\) . It is characterized by

\[
\Phi (x _ {1} \otimes \dots \otimes x _ {m}, y _ {1} \otimes \dots \otimes y _ {m}) = \prod_ {i = 1} ^ {m} \Phi_ {i} (x _ {i}, y _ {i}) \quad (x _ {i} \in \mathrm{E} _ {i}, y _ {i} \in \mathrm{F} _ {i}).\tag{35}
\]

DEFINITION 11. --- Given A-modules \(\mathbf{E}_{i}, \mathbf{F}_{i}\) ( \(i=1,\ldots,m\) ) and, for each \(i\) , a sesquilinear form \(\Phi_{i}\) for J on \(\mathbf{E}_{i} \times \mathbf{F}_{i}\) , the sesquilinear form \(\Phi\) for J on \((\bigotimes_{i}\mathbf{E}_{i}) \times (\bigotimes_{i}\mathbf{F}_{i})\) characterized by (35) is called the tensor product of the sesquilinear forms \(\Phi_{i}\) .

In the case where the \(E_{i}\) and the \(F_{i}\) are equal to a single module E, and where the \(\Phi_{i}\) are equal to a single form \(\Psi\) , one says that \(\Phi\) is the extension of \(\Psi\) to \(\otimes^{m}E\) .

With the notations of def. 11, let us study the maps associated with \(\Phi\) . From formula (24) (no. 6) and (35) we derive the relation

\[
\Phi (x _ {1} \otimes \dots \otimes x _ {m}, y _ {1} \otimes \dots \otimes y _ {m}) = \prod_ {i = 1} ^ {m} \left\langle x _ {i}, d _ {\Phi_ {i}} (y _ {i}) \right\rangle = \prod_ {i = 1} ^ {m} \left\langle y _ {i}, s _ {\Phi_ {i}} (x _ {i}) \right\rangle^ {j}.
\]

We therefore have:

\[
s _ {\Phi} = j _ {s} \circ (s _ {\Phi_ {1}} \otimes \dots \otimes s _ {\Phi_ {m}}), \quad d _ {\Phi} = j _ {d} \circ (d _ {\Phi_ {1}} \otimes \dots \otimes d _ {\Phi_ {m}})\tag{36}
\]

where \(j_s\) (resp. \(j_d\) ) denotes the canonical map of \(\bigotimes_i F_i^*\) into \((\bigotimes_i F_i)^*\) (resp. of \(\bigotimes_i E_i^*\) into \((\bigotimes_i E_i)^*\) ) (chap. III, § 1, nos. 4 and 7).

PROPOSITION 9. — Let A be a commutative field equipped with an automorphism J, E \(_{i}\) , F \(_{i}\) two finite-dimensional vector spaces over A, and Φ \(_{i}\) a sesquilinear form for J on E \(_{i}\) × F \(_{i}\) (1 ≤ i ≤ m). If the forms Φ \(_{i}\) are nondegenerate, then so is their tensor product Φ. In this case the inverse form \(\hat{\Phi}\) of \(\Phi\) is the tensor product of the inverse forms Φ \(_{i}\) .

Indeed, since A is a field, it follows from Props. 6 and 7 of Chap. III, § 1, no. 3 that a tensor product of injective (resp. surjective) linear maps of A-modules is an injective (resp. surjective) linear map. Since the \(s_{\Phi_i}\) are bijective by hypothesis (prop. 6, no. 6), hence so is their tensor product. On the other hand, the canonical map \(j_s\) of \(\bigotimes_i F_i^*\) into \((\bigotimes_i F_i)^*\) is bijective (chap. III, § 1, no. 5, prop. 11). Therefore, by virtue of (36), \(s_\Phi\) is bijective, and this establishes our first assertion (prop. 6, no. 6). Likewise \(d_\Phi\) is bijective.

In the second assertion we have implicitly identified \(\bigotimes_{i} \mathrm{F}_{i}^{*}\) with \((\bigotimes_{i} \mathrm{F}_{i})^{*}\) and \(\bigotimes_{i} \mathrm{E}_{i}^{*}\) with \((\bigotimes_{i} \mathrm{E}_{i})^{*}\) by means of the maps \(j_{s}\) and \(j_{d}\) , which are here isomorphisms. The inverse forms cited in the statement exist since the \(s_{\Phi_{i}}\) , the \(d_{\Phi_{i}}\) , \(s_{\Phi}\) and \(d_{\Phi}\) are bijective (no. 7). Then set \(x' = x_{1}' \otimes \cdots \otimes x_{m}'\) , \(y' = y_{1}' \otimes \cdots \otimes y_{m}'\) ( \(x_{i}' \in \mathrm{E}_{i}^{*}\) , \(y_{i}' \in \mathrm{F}_{i}^{*}\) , \(i = 1, \ldots, m\) ). By definition of the inverse forms, and in view of (36), we have

\[
\begin{array}{l} \widehat {\Phi} (j _ {s} (y ^ {\prime}), j _ {d} (x ^ {\prime})) = \Phi (s _ {\Phi_ {1}} ^ {- 1} (y _ {1} ^ {\prime}) \otimes \dots \otimes s _ {\Phi_ {m}} ^ {- 1} (y _ {m} ^ {\prime}), d _ {\Phi_ {1}} ^ {- 1} (x _ {1} ^ {\prime}) \otimes \dots \otimes d _ {\Phi_ {m}} ^ {- 1} (x _ {m} ^ {\prime})) \\ = \prod_ {i = 1} ^ {m} \Phi_ {i} (s _ {\Phi_ {i}} ^ {- 1} (y _ {i} ^ {\prime}), d _ {\Phi_ {i}} ^ {- 1} (x _ {i} ^ {\prime})) = \prod_ {i = 1} ^ {m} \widehat {\Phi} _ {i} (y _ {i} ^ {\prime}, x _ {i} ^ {\prime}), \end{array}
\]

whence our second assertion.

QED.

Let E and F be two modules over the commutative ring A, and \(\Phi\) a sesquilinear form for J on E × F. The map

\[
(x _ {1}, \dots , x _ {m}, y _ {1}, \dots , y _ {m}) \rightarrow \det (\Phi (x _ {i}, y _ {k})) \quad (x _ {i} \in \mathrm{E}, y _ {i} \in \mathrm{F}, i = 1, \dots , m)
\]

of \(\mathbf{E}^{m}\times(\mathbf{F}^{\mathrm{j}})^{m}\) into A is A-multilinear. It therefore defines a bilinear form \(\Phi'\) on \((\bigotimes^{\overline{m}}\mathbf{E})\times(\bigotimes^{\overline{m}}\mathbf{F}^{\mathrm{j}})\) characterized by

\[
\Phi^ {\prime} (x _ {1} \otimes \dots \otimes x _ {m}, y _ {1} \otimes \dots \otimes y _ {m}) = \det (\Phi (x _ {i}, y _ {k})).
\]

Since the first member is zero when \(x_{i}=x_{k}\) or \(y_{i}=y_{k}\) ( \(i\neq k\) ), \(\Phi^{\prime}\) defines, by passing to quotients, a bilinear form on \((\bigwedge^{m}\mathrm{E})\times(\bigwedge^{m}\mathrm{F}^{\mathrm{J}})\) or again, since \(\bigwedge^{m}\mathrm{F}^{\mathrm{J}}\) is identified with \((\bigwedge^{m}\mathrm{F})^{\mathrm{J}}\) , a form \(\Phi_{(m)}\) sesquilinear for J on \((\bigwedge^{m}\mathrm{E})\times(\bigwedge^{m}\mathrm{F})\) . It is characterized by

\[
\left\{ \begin{array}{l} \Phi_ {(m)} (x _ {1} \wedge \dots \wedge x _ {m}, y _ {1} \wedge \dots \wedge y _ {m}) = \det (\Phi (x _ {i}, y _ {k})) \\ (x _ {i} \in \mathrm{E}, y _ {i} \in \mathrm{F}, i = 1, \dots , m). \end{array} \right.\tag{37}
\]

DEFINITION 12. --- Given two A-modules E, F and a form \(\Phi\) sesquilinear for J on \(E \times F\) , the form \(\Phi_{(m)}\) sesquilinear for J on \((\bigwedge^{m} E) \times (\bigwedge^{m} F)\) characterized by (37) is called the extension of \(\Phi\) to the m-th exterior powers.

With the notation of def. 12, let us study the maps associated with \(\Phi_{(m)}\) . From formula (24) (no. 6) and (37) we derive the relations

\[
\begin{array}{r l} \Phi_ {(m)} (x _ {1} \wedge \dots \wedge x _ {m}, y _ {1} \wedge \dots \wedge y _ {m}) & = \det (\langle x _ {i}, d _ {\Phi} (y _ {k}) \rangle) \\ & = \det (\langle y _ {i}, s _ {\Phi} (x _ {k}) \rangle^ {j}). \end{array}
\]

We therefore have

\[
s _ {\Phi (m)} = k _ {s} \circ (\stackrel {m} {\wedge} s _ {\Phi}), \quad d _ {\Phi_ {m}} = k _ {d} \circ (\stackrel {m} {\wedge} d _ {\Phi}),\tag{38}
\]

where \(k_s\) (resp. \(k_d\) ) denotes the canonical map of \(\bigwedge^m F^*\) into \((\bigwedge^m F)^*\) (resp. of \(\bigwedge^m E^*\) into \((\bigwedge^m E)^*\) ) (cf. Chap. III, § 8, no. 2).

PROPOSITION 10. --- Let A be a commutative field equipped with an automorphism J, E and F two finite-dimensional vector spaces over A, and \(\Phi\) a sesquilinear form for J on E \(\times\) F. If \(\Phi\) is nondegenerate, then its extension \(\Phi_{(m)}\) to the m-th exterior powers is nondegenerate, and the inverse form of \(\Phi_{(m)}\) is the extension to the m-th exterior powers of the inverse form \(\hat{\Phi}\) of \(\Phi\) .

Indeed, as \(s_{\Phi}\) and \(d_{\Phi}\) are bijective by hypothesis (prop. 6, no. 6), and the same holds for their exterior powers (chap. III, § 5, no. 7). On the other hand, the canonical maps \(k_s\) and \(k_d\) are bijective (chap. III, § 8, no. 2, th. 1). Therefore, by virtue of (38), \(s_{\Phi(m)}\) and \(d_{\Phi(m)}\) are bijective, which proves that \(\Phi_{(m)}\) is nondegenerate (prop. 6, no. 6). In the second assertion we have implicitly identified \(\bigwedge^m F^*\) with \((\bigwedge^m F)^*\) and \(\bigwedge^m E^*\) with \((\bigwedge^m E)^*\) by means of the maps \(k_s\) and \(k_d\) , which are here isomorphisms (loc. cit.). The inverse forms considered in the statement exist since \(s_{\Phi}\) , \(d_{\Phi}\) , \(s_{\Phi(m)}\) and \(d_{\Phi(m)}\) are bijective (no. 7). Then set \(x' = x_1' \wedge \ldots \wedge x_m'\) and \(y' = y_1' \wedge \ldots \wedge y_m'\) ( \(x_i' \in E^*\) , \(y_i' \in F^*\) , \(i = 1, \ldots, m\) ). By definition of inverse forms (no. 7) and in view of (38), we have

\[
\begin{array}{r l} \hat {\Phi} _ {(m)} (k _ {s} (y ^ {\prime}), k _ {d} (x ^ {\prime})) & = \Phi_ {(m)} (s _ {\Phi} ^ {- 1} (y _ {1} ^ {\prime}) \wedge \dots \wedge s _ {\Phi} ^ {- 1} (y _ {m} ^ {\prime}), d _ {\Phi} ^ {- 1} (x _ {1} ^ {\prime}) \wedge \dots \wedge d _ {\Phi} ^ {- 1} (x _ {m} ^ {\prime})) \\ & = \det (\Phi (s _ {\Phi} ^ {- 1} (y _ {i} ^ {\prime}), d _ {\Phi} ^ {- 1} (x _ {k} ^ {\prime}))) = \det (\hat {\Phi} (y _ {i} ^ {\prime}, x _ {k} ^ {\prime})) \end{array}
\]

whence our second assertion.

Remark. --- Let E be a free A-module, and let θ be the canonical isomorphism from ∧E onto the submodule of antisymmetrized tensors of order m (chap. III, § 5, no. 6, prop. 6). Let Φ be a sesquilinear form on E, let Φ(m) be the extension of Φ to ∧E, and let Θ be the sesquilinear form on ∧E which is the inverse image with respect to θ of the extension of Φ to ⊗E. By the definition of θ and of the antisymmetrization of a tensor, and by (35), we have

\[
\Theta (x _ {1} \wedge \dots \wedge x _ {m}, y _ {1} \wedge \dots \wedge y _ {m}) = \sum_ {\sigma , \tau} \varepsilon_ {\sigma} \varepsilon_ {\tau} \Phi (x _ {\sigma (1)}, y _ {\tau (1)}) \dots \Phi (x _ {\sigma (m)}, y _ {\tau (m)})
\]

where \(\sigma\) and \(\tau\) range over the symmetric group \(\mathfrak{S}_{m}\) . By the formula for computing determinants and formula (37), this expression can be written

\[
\sum_ {\tau \in \mathfrak {S} _ {m}} \varepsilon_ {\tau} \det (\Phi (x _ {i}, y _ {\tau (k)})) = m! \det (\Phi (x _ {i}, y _ {k}));
\]

In other words, we have \(\Theta = m!\Phi_{(m)}\) .

\subsection{Matrix calculations.}

In this no., we propose to relax the matrix calculus introduced in Chap. II, § 6, and to apply it to translate certain results proved in this section.

I. --- Let I and K be two finite index sets, H a nonempty set, and \(M = (m_{ik})_{(i,k) \in \mathbb{I} \times \mathbb{K}}\) a matrix over H (chap. II, § 6, no. 1, def. 1).

The transpose of \(M\) is called, and denoted \({}^t M\) , the matrix \((m_{ki}')_{(k,i) \in \mathbb{K} \times \mathbf{I}}\) satisfying \(m_{ki}' = m_{ik} ((i, k) \in \mathbf{I} \times \mathbf{K})\) . We obviously have

\[
^ {t} (^ {t} M) = M.\tag{39}
\]

This generalizes the notion introduced in Chap. II, § 6, no. 6.

Suppose that H is a commutative group (written additively). The set of matrices over H with I and K as index sets admits a structure of commutative group, since it is the set of maps from I × K to H. This group is written additively.

Let H', H'' be two nonempty sets, H an abelian group (written additively), and \(f:(h',h'')\to h'h''\) a map of \(\mathbf{H}'\times\mathbf{H}''\) into H. Given two matrices

\[
M ^ {\prime} = (m _ {i k} ^ {\prime}) _ {(i, k) \in \mathrm{I} \times \mathrm{K}}, \quad M ^ {\prime \prime} = (m _ {k l} ^ {\prime \prime}) _ {(k, l) \in \mathrm{K} \times \mathrm{L}}
\]

on H' and H'' respectively, such that the set K of column indices of M' is equal to the set of row indices of M'' , one calls the product of M' and M'' (following f) and denotes it M'M'' the matrix

\[
M ^ {\prime}. M ^ {\prime \prime} = (\sum_ {k \in \mathbf {K}} m _ {i k} ^ {\prime} m _ {k l} ^ {\prime \prime}) _ {(i, l) \in \mathbf {I} \times \mathbf {L}}\tag{40}
\]

on H. This generalizes the notion introduced in chap. II, § 6, no. 4. If \(\mathbf{H}' = \mathbf{H}'' = \mathbf{H}\) and if H is a ring, the product \(M'M''\) will, unless expressly stated otherwise, be computed "in H", that is, according to the map \((x, y) \to xy\) . When \(\mathbf{H}'\) and \(\mathbf{H}''\) are commutative groups (written additively) and \(f\) is bilinear, one has

\[
\left\{ \begin{array}{l} (M ^ {\prime} + M _ {1} ^ {\prime}) M ^ {\prime \prime} = M ^ {\prime} M ^ {\prime \prime} + M _ {1} ^ {\prime} M ^ {\prime \prime}, \\ M ^ {\prime} (M ^ {\prime \prime} + M _ {1} ^ {\prime \prime}) = M ^ {\prime} M ^ {\prime \prime} + M ^ {\prime} M _ {1} ^ {\prime \prime}, \end{array} \right.\tag{41}
\]

where \(M'\) , \(M_1'\) are matrices over \(H'\) , \(M''\) , \(M_1''\) matrices over \(H''\) , and where the sums and products written are assumed to be defined. Let \(M'\) , \(M''\) be matrices over the sets \(H'\) , \(H''\) , and \(f^0\) the map of \(H'' \times H'\) into H defined by \((h'', h') \to h' h''\) ; then one has

\[
{ } ^ { t } ( M ^ { \prime } M ^ { \prime \prime } ) = { } ^ { t } M ^ { \prime \prime } . { } ^ { t } M ^ { \prime }\tag{42}
\]

where the product in the first (resp. second) member is computed according to \(f\) (resp. \(f^{0}\) ).

In the case where \(H' = H'' = H\) is a ring, one recovers formula (12) of chap. II, §6, no. 6.

Let A be a ring, J an antiautomorphism of A. For every matrix \(M = (m_{ik})\) over A, we shall denote by \(M^{J}\) the matrix \((m_{ik}^{J})\) . Let \(M_{1}\) , \(M_{2}\) be two matrices over A such that \(M_{1}M_{2}\) is defined. Since J is an isomorphism of A onto the opposite ring \(A^{0}\) , one has \((M_{1}M_{2})^{J} = M_{1}^{J}\) . \(M_{2}^{J}\) where the first (resp. second) member is computed in A (resp. \(A^{0}\) ). In view of (42) and (39), this gives

\[
(M _ {1} M _ {2}) ^ {\mathrm{J}} = ^ {t} (^ {t} M _ {2} ^ {\mathrm{J}}. ^ {t} M _ {1} ^ {\mathrm{J}})\tag{43}
\]

where both members are computed in A.

Let \(H_{1}\) , \(H_{2}\) , \(H_{3}\) , \(H_{12}\) , \(H_{23}\) and H be commutative groups (written additively), \(f_{12}: H_{1} \times H_{2} \to H_{12}\) , \(f_{23}: H_{2} \times H_{3} \to H_{23}\) , \(f_{3}: H_{12} \times H_{3} \to H\) , \(f_{1}: H_{1} \times H_{23} \to H\) maps, and let \(M_{1}\) , \(M_{2}\) , \(M_{3}\) be matrices over \(H_{1}\) , \(H_{2}\) , \(H_{3}\) respectively. If \(f_{3}(f_{12}(x_{1}, x_{2}), x_{3}) = f_{1}(x_{1}, f_{23}(x_{2}, x_{3}))\) for all \(x_{i} \in H_{i}\) (i = 1, 2, 3), then the products \((M_{1}M_{2})M_{3}\) and \(M_{1}(M_{2}M_{3})\) (computed according to \(f_{12}\) , \(f_{3}\) , \(f_{23}\) and \(f_{1}\) ), if they are defined, are equal; they will be denoted \(M_{1}M_{2}M_{3}\) . When \(H_{1} = H_{2} = H_{3} = H_{12} = H_{23} = H\) , that H is a ring, and that \(f_{12}\) , \(f_{23}\) , \(f_{3}\) , \(f_{1}\) are equal to the map \((x, y) \to xy\) , the preceding condition expresses the associativity of the latter, and is therefore satisfied. Analogous conventions will be made for products of more than three factors.

Let A, B be two rings, \(M = (m_{ik})_{(i,k) \in I \times K}\) and \(M' = (m_{ik}')_{(i,k) \in I \times K}\) be two matrices over an (A, B)-bimodule G (no. 1). If, for every one-row matrix \(L = (a_i)_{i \in I}\) with elements in A and every one-column matrix \(C = (b_k)_{k \in K}\) with elements in B, one has \(L.M.C = L.M'.C\) (the products being computed according to the maps defining the structure of (A, B)-bimodule of G), then the matrices \(M\) and \(M'\) are equal. Indeed, if one takes \(a_i = 1\) , \(a_s = 0\) for \(s \neq i\) , \(b_k = 1\) , \(b_t = 0\) for \(t \neq k\) , the matrices \(L.M.C\) and \(L.M'.C\) , which are scalar matrices, are respectively equal to \(m_{ik}\) and \(m_{ik}'\) .

\begin{enumerate}
\def\labelenumi{\Roman{enumi}.}
\setcounter{enumi}{1}
\tightlist
\item
  Consider a ring A and a (right or left) A-module E, admitting a finite basis \((e_{i})_{i\in I}\) . For every element x of E, the one-column matrix formed by the components \(x_{i}\) ( \(i\in I\) ) of x with respect to \((e_{i})\) is called the matrix of x with respect to the basis \((e_{i})\) , and is denoted \(M(x)\) or x (cf. chap. II, § 6, no. 4); in the calculations it will be convenient, in order to recall that the index i is a row index, to adjoin to it a column index capable of taking only one value, and to write \((x_{i0})\) for the matrix \(M(x)\) .
\end{enumerate}

Now consider two A-modules (left or right) E and F, having finite bases \((e_{i})_{i\in I}\) and \((f_{k})_{k\in K}\) respectively; let \((f_{k}^{*})\) be the basis of \(F^{*}\) dual to \((f_{k})\) . We shall define the matrix, with respect to these bases, of a map u from E to F in the following four cases:

(D) E and F are right modules, and u is A-linear;

(G) E and F are left modules, and u is A-linear;

(GD) E is a left module, F a right module, A is equipped with an antiautomorphism J, u is Z-linear and satisfies \(u(ax) = u(x)a^{J}\) ( \(a \in A, x \in E\) ) (in other words, u is an A-linear map from \(E^{J}\) to F (no. 2, def. 5)).

(DG) E is a right A-module, F a left A-module, A is equipped with an antiautomorphism J, u is Z-linear and satisfies \(u(xa) = a^{J}u(x)\) ( \(x \in E, a \in A\) ) (in other words, u is an A-linear map from \(E^{J}\) into F).

In each of these four cases, the matrix of the map u is, by definition, the matrix \((u_{ki})_{(k,i)\in\mathbb{K}\times\mathbb{I}}\) such that

\[
u _ {k i} = \left\langle u (e _ {i}), f _ {k} ^ {*} \right\rangle .\tag{44}
\]

This definition coincides, in case (D), with that given in Chap. II, § 6, no. 3. Under these conditions, the matrix \(M(u(x))\) of the image of an element \(x\) of E is given by the following formulas:

\[
M (u (x)) = M (u). M (x)\tag{45 D}
\]

\[
{ } ^ { t } M ( u ( x ) ) = { } ^ { t } M ( x ) . { } ^ { t } M ( u )\tag{45 G}
\]

\[
M (u (x)) = M (u). M (x) ^ {\mathrm{J}}\tag{45 GD}
\]

\[
{ } ^ { t } M ( u ( x ) ) = { } ^ { t } M ( x ) ^ { j } . { } ^ { t } M ( u ) .\tag{45 DG}
\]

Let us verify, for example, (45 DG), the other verifications being analogous and slightly easier. Set \(x = \sum_{i} e_i x_{io}, u(x) = \sum_{k} y_{ko} f_k\) ; we have \(u(x) = u(\sum_{i} e_i x_{io}) = \sum_{i} x_{io}^{\mathbf{J}} u(e_i) = \sum_{i,k} x_{io}^{\mathbf{J}} u_{ki} f_k\) ; whence \(y_{ko} = \sum_{i} x_{io}^{\mathbf{J}} u_{ki}\) ; in order to place the two indices \(i\) next to each other, consider the transposed matrices \(^t M(x) = (x_{oi}')\) where \(x_{oi}' = x_{io}\) , and \(^t M(u) = (u_{ik}')\) where \(u_{ik}' = u_{ki}\) ; we then have \(y_{ko} = \sum_{i} x_{oi}' u_{ik}'\) ; since the right-hand side is the element with index \(k\) of the one-row matrix \(^t M(x)^{\mathbf{J}}\) . \(^t M(u)\) , formula (45 DG) is verified.

Remarks. -- 1) When A is commutative, (45 G) reduces to (45 D), and (45 DG) to (45 GD), by means of the formula \(t(M'M'') = 'M'''.'M'\) (cf. (42)), where both sides are here computed in A. 2) Let E, F, G be three left modules having finite bases, and \(u: \mathrm{E} \to \mathrm{F}\) , \(\nu: \mathrm{F} \to \mathrm{G}\) A-linear maps. It follows from (45 G) that one has

\[
{ } ^ { t } M ( \nu \circ u ) = { } ^ { t } M ( u ) . { } ^ { t } M ( \nu ) .\tag{46}
\]

Indeed, we have, for any \(x \in E\) ,

\[
\begin{array}{r l} ^ {t} M (x). ^ {t} M (\nu \circ u) & = ^ {t} M (\nu (u (x))) = ^ {t} M (u (x)). ^ {t} M (\nu) \\ & = ^ {t} M (x). ^ {t} M (u). ^ {t} M (\nu), \end{array}
\]

whence (46).

Recall that, in the case of right modules, we have

\[
M (\nu \circ u) = M (\nu) M (u).
\]

\begin{enumerate}
\def\labelenumi{\Roman{enumi}.}
\setcounter{enumi}{2}
\tightlist
\item
  From now on, we denote by A a ring, by B a ring (resp. a ring equipped with an antiautomorphism J, for which we set \(J' = J^{-1}\) ), by E a left A-module having a finite basis \((e_i)_{i \in I}\) and by F a right (resp. left) B-module having a finite basis \((f_k)_{k \in K}\) . We denote by \((e_i^*)\) and \((f_k^*)\) the dual bases of \(E^*\) and \(F^*\) . Unless expressly stated otherwise, the matrices considered are taken with respect to these bases.
\end{enumerate}

Let G be an (A, B)-bimodule (no. 1), Φ a bilinear (resp. sesquilinear on the right for J) map of E × F into G, and \(R = (\Phi(e_i, f_k))\) the matrix of Φ. Then, for \(x \in E\) and \(y \in F\) , formula (6) of no. 1 (resp. (8) of no. 2) is written, by virtue of the above conventions,

\[
\Phi (x, y) = ^ {t} M (x). R. M (y) \quad (\text { resp. } \Phi (x, y) = ^ {t} M (x). R. M (y) ^ {\mathrm{J}}),\tag{47}
\]

where the products are computed according to the maps which define the structure of the (A, B)-bimodule G; in particular, if A = B = G (in which case Φ is a form), the products are computed in A.

Let E' be a left A-module having a finite basis \((e_{s}^{\prime})_{s\in S}\) , F' a right (resp. left) A-module having a finite basis \((f_{t}^{\prime})_{t\in T}\) , \(u:E\to E^{\prime}\) and \(\nu:F\to F^{\prime}\) A-linear maps, and \(\Phi^{\prime}\) a bilinear map (resp. sesquilinear on the right for J) from \(E^{\prime}\times F^{\prime}\) into G. Denote by \(\Phi\) the inverse image of \(\Phi^{\prime}\) (relative to u and \(\nu\) ), U, V, R, \(R^{\prime}\) the matrices of u, \(\nu\) , \(\Phi\) , \(\Phi^{\prime}\) with respect to the bases considered. One then has

\[
R = ^ {t} U. R ^ {\prime}. V \quad (\text { resp. } R = ^ {t} U. R ^ {\prime}. V ^ {J}),\tag{48}
\]

the products being computed as in (47). Indeed, whatever \(x \in \mathbf{E}\) and \(y \in \mathbf{F}\) , we have by definition \(\Phi(x, y) = \Phi'(u(x), v(y))\) , whence, by (47),

\[
\begin{array}{c} ^ {t} M (x). R. M (y) = ^ {t} M (u (x)). R ^ {\prime}. M (\nu (y)) \\ (\text {resp.} ^ {t} M (x). R. M (y) ^ {\mathrm{J}} = ^ {t} M (u (x)). R ^ {\prime}. M (\nu (y)) ^ {\mathrm{J}}); \end{array}
\]

by (45 G) and (45 D) (resp. (45 G)) and (43) one deduces

\[
\begin{array}{r l} ^ {t} M (x). R. M (y) & = ^ {t} M (x). ^ {t} U. R ^ {\prime}. V. M (y) \\ (\text {resp.} ^ {t} M (x). R. M (y) ^ {J} & = ^ {t} M (x). ^ {t} U. R ^ {\prime}. ^ {t} (^ {t} M (y). ^ {t} V) ^ {J} \\ & = ^ {t} M (x). ^ {t} U. R ^ {\prime}. V ^ {J}. M (y) ^ {J}); \end{array}
\]

this proves our assertion.

\begin{enumerate}
\def\labelenumi{\Roman{enumi}.}
\setcounter{enumi}{3}
\tightlist
\item
  --- We assume here that the rings A and B are equal, and we denote by \(\Phi\) a bilinear (resp. sesquilinear on the right for J) form on E × F, and by \(R\) its matrix. Let us compute the matrices of the maps \(s_{\Phi}\) and \(d_{\Phi}\) associated with \(\Phi\) , which we shall denote by \(s\) and \(d\) to simplify. Since one has \(\Phi(x, y) = \langle y, s(x) \rangle = \langle x, d(y) \rangle\) by (23), no. 6 (resp. \(\Phi(x, y) = \langle x, d(y) \rangle = \langle y, s(x) \rangle^{J}\) by (24), no. 6), one has \(\Phi(e_i, f_k) = \langle f_k, s(e_i) \rangle = \langle e_i, d(f_k) \rangle\) (resp. \(\Phi(e_i, f_k) = \langle e_i, d(f_k) \rangle = \langle f_k, s(e_i) \rangle^{J}\) ), whence, by (44) and since ( \(e_i\) ) is the dual basis of ( \(e_i^*\) ) and ( \(f_k\) ) the dual basis of ( \(f_k^*\) ):
\end{enumerate}

\[
M (d) = R,   M (s) = ^ {t} R \quad (\text { resp. }   M (d) = R,   M (s) = ^ {t} R ^ {J ^ {\prime}}).\tag{49}
\]

Remarks. --- 1) When A is a field, the linear maps s and d have the same rank. We see here that their matrices \(M(s)\) and \(M(d)\) have the same rank; indeed, a matrix over A and its transpose have the same rank (chap. II, § 6, no. 7, prop. 3) and, when \(\Phi\) is sesquilinear, the equality of the ranks of \(R\) over A and of \(^tR\) over A \(^0\) (ibid.) and the fact that J' is an isomorphism of A \(^0\) onto A, imply the equality of the ranks of \(R\) and of \(^tR^{J'}\) over A.

\begin{enumerate}
\def\labelenumi{\arabic{enumi})}
\setcounter{enumi}{1}
\tightlist
\item
  If M and N are right A-modules having finite bases \((m_{i})\) and \((n_{k})\) , \(\Phi\) a sesquilinear form on the left for J on \(M \times N\) (no. 6, Remark), s and d its associated maps, and \(R = (\Phi(m_{i}, n_{k}))\) its matrix, formulas (26) of no. 6 show that one has
\end{enumerate}

\[
M (d) = R ^ {J ^ {\prime}}, \quad M (s) = ^ {t} R.
\]

Suppose now that the maps s and d associated with \(\Phi\) are bijective and let us compute the matrix \(\hat{R}\) of the inverse form of \(\Phi\) (no. 7). When \(\Phi\) is bilinear, \(\Phi\) is the inverse image of \(\hat{\Phi}\) with respect to the linear maps \(s: \mathrm{E} \to \mathrm{F}^*\) and \(d: \mathrm{F} \to \mathrm{E}^*\) ; one therefore has, by virtue of (48) and (49), \(R = R\) . \(\hat{R}\) . \(R\) , whence, since \(R\) is invertible ( \(d\) being bijective), \(\hat{R} = R^{-1}\) . This formula extends to the case where \(\Phi\) is sesquilinear, for, if one considers \(\Phi\) as a bilinear form on \(\mathrm{E} \times \mathrm{F}^{\mathrm{j}}\) , and if one identifies \((\mathrm{F}^{\mathrm{j}})^{*}\) with \((\mathrm{F}^{*})^{\mathrm{j}}\) (cf. no. 9), the inverse form of this bilinear form coincides with \(\hat{\Phi}\) considered as a bilinear form on \((\mathrm{F}^{*})^{\mathrm{j}} \times \mathrm{E}^{*}\) . In both cases the matrix of the inverse form of \(\Phi\) is the inverse of the matrix of \(\Phi\) .

Finally, let E' be a left A-module, F' a right (resp. left) A-module, both admitting finite bases \((e_{s}^{\prime})\) and \((f_{t}^{\prime})\) ; let \(\Phi^{\prime}\) be a bilinear form (resp. sesquilinear form with respect to J) on \(E^{\prime} \times F^{\prime}\) , and let \(R^{\prime}\) be its matrix. Suppose \(s_{\Phi}\) and \(d_{\Phi}\) bijective. Let \(u : E \to E^{\prime}\) and \(\nu : F \to F^{\prime}\) be linear maps, \(u^{*} : F^{\prime} \to F\) and \(\upsilon^{*} : E^{\prime} \to E\) their adjoints (no. 8, prop. 7); denote by U, V, \(U^{*}\) , \(V^{*}\) the matrices of u, \(\upsilon\) , \(u^{*}\) , \(\upsilon^{*}\) with respect to the given bases. One then has

\[
U ^ {*} = R ^ {- 1}. ^ {t} U. R ^ {\prime}, \quad {} ^ {t} V ^ {*} = R ^ {\prime}. V. R ^ {- 1}\tag{50}
\]

\[
(\text { resp. } U ^ {* \mathrm{J}} = R ^ {- 1}. ^ {t} U. R ^ {\prime}, ^ {t} V ^ {*} = R ^ {\prime}. V ^ {\mathrm{J}}. R ^ {- 1}).
\]

Indeed, whatever \(x \in E\) and \(y \in F'\) , one has \(\Phi'(u(x), y) = \Phi(x, u^*(y))\) (no. 8, def. 10). Hence, when \(\Phi\) is bilinear, by virtue of (47), \(^t M(u(x)). R'. M(y) = ^t M(x). R. M(u^*(y))\) ; this gives, by virtue of (45 G) and (45 D), \(^t M(x). ^t U. R'. M(y) = ^t M(x). R. U^*. M(y)\) , whence \(^t U. R' = R. U^*\) and the first announced formula, since, \(d\) being bijective, \(R\) is invertible. When \(\Phi\) is sesquilinear (47) and (45 G) give \(^t M(x). ^t U. R'. M(y)^J = ^t M(x). R. (^{t}M(y). ^{t}U^*)^{J}\) ; now, according to (43), we have \((^{t}M(y). ^{t}U^*)^{J} = ^{t}(^{u}U^{*J}. ^{u}M(y)^{J})\) , whence \((^{t}M(y). ^{t}U^*)^{J} = U^{*J}. M(y)^{J}\) ; it follows therefore \(^t M(x). ^{t}U. R'. M(y)^{J} = ^{t}M(x). R. U^{*J}. M(y)^{J}\) , whence \(^t U. R' = R. U^{*J}\) , and \(U^{*J} = R^{-1}. ^{t}U. R'\) . The verification of the formulas for \(V^*\) is analogous.

Exercises. --- 1) Let A be a commutative field, E a vector space over A admitting an infinite countable basis \((e_n)_{n \geqslant 1}\) . We define a bilinear form \(\Phi\) on E by setting \(\Phi(e_{i+1}, e_i) = 1\) for \(i \geqslant 1\) , \(\Phi(e_k, e_j) = 0\) for \(k \neq j + 1\) and \(j \geqslant 1\) . Show that the linear map \(d_\Phi\) associated on the right with \(\Phi\) is injective, but the linear map \(s_\Phi\) associated on the left with \(\Phi\) is not injective.

\begin{enumerate}
\def\labelenumi{\arabic{enumi})}
\setcounter{enumi}{1}
\item
  Let E be the Z-module direct sum of Z and Z/(2), and let E* be its dual (isomorphic to Z). Show that the bilinear form \((x, x') \to \langle x, x' \rangle\) on \(E \times E^*\) is such that the associated linear map on the right is injective, but the associated linear map on the left is not.
\item
  Give an example of a bilinear form \(\Phi\) defined on a product \(\mathbf{E} \times \mathbf{F}\) of two vector spaces, such that \(d_{\Phi}\) is bijective, \(s_{\Phi}\) injective but not bijective (take \(\mathbf{E}\) of infinite dimension and F equal to the dual \(\mathbf{E}^*\) of \(\mathbf{E}\) ; cf. chap. II, § 5, exerc. 3).
\item
  Let A be a ring equipped with an antiautomorphism J, let E be a left A-module, and let G be an (A, A)-bimodule, and \(\Phi\) a map from \(\mathrm{E} \times \mathrm{E}\) into G, sesquilinear on the right for J. Prove the identity (where \(\mathrm{Q}(x) = \Phi(x, x)\) ):
\end{enumerate}

\[
\begin{array}{r l} & {2 \Phi (x, y) (\mu^ {\mathrm{J}} \lambda^ {\mathrm{J}} - \lambda^ {\mathrm{J}} \mu^ {\mathrm{J}}) = \mathrm{Q} (x - \mu \lambda y) - \mathrm{Q} (x + \mu \lambda y) + \mu \mathrm{Q} (x + \lambda y)} \\ & {- \mu \mathrm{Q} (x - \lambda y) + \mathrm{Q} (x + \mu y) \lambda^ {\mathrm{J}} - \mathrm{Q} (x - \mu y) \lambda^ {\mathrm{J}} + \mu \mathrm{Q} (x - y) \lambda^ {\mathrm{J}} - \mu \mathrm{Q} (x + y) \lambda^ {\mathrm{J}}.} \end{array}
\]

\begin{enumerate}
\def\labelenumi{\arabic{enumi})}
\setcounter{enumi}{4}
\tightlist
\item
  Let K be a commutative field of characteristic 2, A a separable quadratic extension of K; we have \(\mathrm{A} = \mathrm{K}(\theta)\) , where \(\theta\) is a root of an irreducible polynomial \(\mathrm{X}^2 +\mathrm{X} + \beta\) of K{[}X{]} and the K-automorphism J of A, distinct from the identity, is such that \(\theta^{j} = \theta +1\) (chap. V, § 11, exerc. 8). Show that if E and G are vector spaces over A, \(\Phi\) a sesquilinear map (for J) from \(\mathrm{E}\times \mathrm{E}\) into G, we have, setting \(\mathrm{Q}(x) = \Phi (x,x)\) ,
\end{enumerate}

\[
\Phi (x, y) = \mathrm{Q} (\theta x + y) - \beta \mathrm{Q} (x) - \mathrm{Q} (y) - (\theta + 1) (\mathrm{Q} (x + y) - \mathrm{Q} (x) - \mathrm{Q} (y)).
\]

\begin{enumerate}
\def\labelenumi{\arabic{enumi})}
\setcounter{enumi}{5}
\tightlist
\item
  Let A be a field, E a vector space over A, \(\Phi\) a sesquilinear form on E, \(u\) an endomorphism of E.
\end{enumerate}

\begin{enumerate}
\def\labelenumi{\alph{enumi})}
\item
  For there to exist one and only one endomorphism \(u^*\) of E such that \(\Phi(u(x), y) = \Phi(x, u^*(y))\) for \(x, y\) in E, it is necessary and sufficient that \(d_\Phi\) be injective and that \(^t u(d_\Phi(E)) \subset d_\Phi(E)\) .
\item
  Give an example where E is infinite-dimensional and \(d_{\Phi}\) injective, but where \(u(d_{\Phi}(E))\) is not contained in \(d_{\Phi}(E)\) .
\end{enumerate}

\begin{enumerate}
\def\labelenumi{\arabic{enumi})}
\setcounter{enumi}{6}
\tightlist
\item
  Let E, \(E_{1}\) be two A-modules, \(\Phi\) (resp. \(\Phi_{1}\) ) a sesquilinear form on E (resp. \(E_{1}\) ). Assume that \(\Phi_{1}\) is nondegenerate and that there exists an element \(\alpha\in A\) and a bijection u from E onto \(E_{1}\) such that \(\Phi_{1}(u(x), u(y))=\Phi(x,y)\alpha\) for all x, y in E. Show that: \(1^{\circ}\Phi\) is nondegenerate; \(2^{\circ}u\) is linear; \(3^{\circ}\) if \(E_{1}\) is a faithful A-module, so is E, and \(\alpha\) is not a right zero divisor in A; \(4^{\circ}\) if \(\Phi_{1}\) takes values in A that are not left zero divisors, so is \(\Phi\) .
\end{enumerate}

¶ 8) Let A be a field, E₁, E₂ two nonzero vector spaces over A, Φ₁ (resp. Φ₂) a nondegenerate sesquilinear form on E₁ (resp. E₂) for an antiautomorphism J₁ (resp. J₂) of A. Let u be a linear map from E₁ onto E₂ such that the relation Φ₁(x, y) = 0 implies Φ₂(u(x), u(y)) = 0.

\begin{enumerate}
\def\labelenumi{\alph{enumi})}
\item
  Show that \(u\) is a bijection from \(\mathbf{E}_{1}\) onto \(\mathbf{E}_{2}\) . (If \(u(0)\) were not reduced to 0, show that there would exist in \(\mathbf{E}_{1}\) two vectors \(a\) , \(b\) such that \(u(a) \neq 0\) , \(u(b) = 0\) and \(\Phi_{1}(a, b) \neq 0\) ; if H is the hyperplane of \(x \in \mathbf{E}_{1}\) such that \(\Phi_{1}(a, x) = 0\) , note that one would have \(u(\mathrm{H}) = \mathrm{E}_{2}\) ).
\item
  Show that if dim \(E_{1} \geqslant 2\) , there exists \(\alpha \in A\) such that one has \(\Phi_{2}(u(x), u(y)) = \Phi_{1}(x, y) \alpha\) for all x, y in \(E_{1}\) . (For every \(y \in E_{1}\) , show that there exists an element \(m(y) \in A\) such that \(\Phi_{2}(u(x), u(y)) = \Phi_{1}(x, y)m(y)\) for all \(x \in E_{1}\) , and that if y and \(y'\) are linearly independent in \(E_{1}\) , one has \(m(y + y') = m(y) = m(y')\) ).
\end{enumerate}

\begin{enumerate}
\def\labelenumi{\arabic{enumi})}
\setcounter{enumi}{8}
\tightlist
\item
  Let A be a field, E, F be two left vector spaces over A, \(\Phi\) a nondegenerate sesquilinear form on \(E \times F\) for an antiautomorphism J of A.
\end{enumerate}

\begin{enumerate}
\def\labelenumi{\alph{enumi})}
\item
  Let M be a subspace of E, and N a subspace of F such that \(N \supset M^{0}\) and \(M \supset N^{0}\) . Show that if one of the spaces \(N/M^{0}\) , \(M/N^{0}\) is finite-dimensional, then so is the other, and the dimensions of these two spaces are equal.
\item
  Let M, \(M'\) be two subspaces of E such that \(M^{00} = M\) and \(M'\) be of finite dimension; show that we have \((M \cap M')^{0} = M^{0} + M'^{0}\) and \((M + M')^{00} = M + M'\) . (Applying a) to the subspaces \(M'\) and \(M^{0} + M'^{0}\) , show that \(\dim(M \cap M') = \text{codim}(M^{0} + M'^{0})\) ; applying a) to the subspaces \(M + M'\) and \(M^{0}\) , show that
\end{enumerate}

\[
\dim ((M + M ^ {\prime}) ^ {0 0} / M) = \dim ((M + M ^ {\prime}) / M).
\]

\begin{enumerate}
\def\labelenumi{\alph{enumi})}
\setcounter{enumi}{2}
\item
  If E = F and if M is a subspace of E such that \(E = M^{0} + M^{00}\) , show that E is the direct sum of \(M^{0}\) and \(M^{00}\) .
\item
  Let E be a vector space over a commutative field A admitting an infinite countable basis \((e_{n})_{n\geqslant0}\) , and let \(\Phi\) be the symmetric bilinear form on E such that \(\Phi(e_{n}, e_{n}) = 1\) for all n, \(\Phi(e_{i}, e_{j}) = 0\) for \(i \geqslant 1\) , \(j \geqslant 1\) and \(i \neq j\) , and \(\Phi(e_{0}, e_{n}) = 1\) for all \(n \geqslant 1\) . Show that \(\Phi\) is nondegenerate. Let M (resp. N) be the subspace of E spanned by the \(e_{2k}\) (resp. \(e_{2k-1}\) ) for \(k \geqslant 1\) , and let H = M + N, which is a hyperplane in E. Show that we have \(M^{0} = N\) , \(N^{0} = M\) , \(H^{00} = E \neq H\) , \((M \cap N)^{0} \neq M^{0} + N^{0}\) and \((M + N)^{00} \neq M + N\) , although \(M^{00} = M\) , \(N^{00} = N\) ; if L is the 2-dimensional subspace spanned by \(e_{0}\) and \(e_{1}\) , one has \((L \cap H)^{0} \neq L^{0} + H^{0}\) .
\end{enumerate}

¶ 10) Let E, E' be two left vector spaces over fields A, A' respectively, of dimension ≥3; let ℱ(E) (resp. ℱ(E')) be the lattice-ordered set (for the inclusion relation) formed by the finite-dimensional subspaces of E (resp. E').

\begin{enumerate}
\def\labelenumi{\alph{enumi})}
\item
  Let \(p\) be a map from \(\mathfrak{F}(\mathrm{E})\) into \(\mathfrak{F}(\mathrm{E}')\) such that for every \(\mathbf{M} \in \mathfrak{F}(\mathrm{E})\) , \(\dim p(\mathbf{M}) = \dim \mathbf{M}\) , and such that for every pair \((\mathbf{M}, \mathbf{N})\) of elements of \(\mathfrak{F}(\mathrm{E})\) , \(p(\mathbf{M} + \mathbf{N}) = p(\mathbf{M}) + p(\mathbf{N})\) . Show that \(p\) is injective; if \(p\) is bijective, there exists a bijective semilinear map \(u\) of \(\mathrm{E}\) into \(\mathrm{E}'\) such that one has \(u(\mathbf{M}) = p(\mathbf{M})\) for all \(\mathbf{M} \in \mathfrak{F}(\mathrm{E})\) (use exercise 10 of chap. II, 2nd ed., App. III).
\item
  Give an example where \(A' = A\) is commutative, \(E' = E\) is of finite dimension, and where there exists a map p from \(\mathfrak{F}(E)\) into itself, such that \(\dim p(M) = \dim M\) , \(p(M + N) = p(M) + p(N)\) , \(p(M \cap N) = p(M) \cap p(N)\) , but there exists no injective semilinear map u from E into itself such that \(u(M) = p(M)\) for \(M \in \mathfrak{F}(E)\) . (Consider the case where there exists an overfield \(A''\) of A of finite degree and isomorphic to A, for example \(\mathbf{A} = \mathbf{F}_{p}(\mathbf{X})\) , where \(p\) is prime; \(\mathrm{E}'' = \mathrm{A}'' \otimes_{\mathbf{A}} \mathrm{E}\) is then a vector space of the same dimension over \(\mathrm{A}''\) as \(\mathrm{E}\) over \(\mathrm{A}\) ; consider the map \(\mathrm{M} \to \mathrm{A}'' \otimes_{\mathbf{A}} \mathrm{M}\) .
\item
  We assume \(A' = A\) . Let \(\omega\) be a bijective map from \(\mathfrak{F}(E)\) onto the set \(\mathfrak{F}'(E')\) of subspaces of finite codimension of \(E'\) , such that \(\text{codim } \omega(M) = \dim M\) , and \(\omega(M + N) = \omega(M) \cap \omega(N)\) . Show that there exists a sesquilinear form \(\Phi\) on \(E \times E'\) , nondegenerate on the right, and such that \(\omega(M) = M^0\) for all \(M \in \mathfrak{F}(E)\) . (Use th. 1 of chap. II, § 4, no. 6).
\end{enumerate}

¶ 11) a) Let A be a ring artinian on the left and on the right (chap. VIII, § 2, no. 3). Show that the following conditions are equivalent: 1° the right annihilator (resp. left annihilator) of every left ideal (resp. right ideal) ≠ A is not reduced to 0; 2° the dual of every simple left A-module (resp. right A-module) is not reduced to 0; 3° the dual of every finitely generated left A-module (resp. right A-module) is not reduced to 0. A ring A is said to satisfy condition (N \(_{s}\) ) (resp. (N \(_{a}\) )) if it satisfies these conditions for left A-modules (resp. right A-modules).

\begin{enumerate}
\def\labelenumi{\alph{enumi})}
\setcounter{enumi}{1}
\item
  Let A be a ring artinian on the left and on the right satisfying condition (N \(_{a}\) ), E (resp. F) a free left A-module (resp. right A-module) having a countable basis over A, Φ a bilinear form on E × F such that s \(_{\Phi}\) is injective. Let M be a free submodule of E; show that there exists a free submodule N of F and a basis (e \(_{n}\) ) (resp. (f \(_{n}\) )) of M (resp. N) such that Φ(e \(_{i}\) , f \(_{j}\) ) = δ \(_{ij}\) ; moreover one can take for e \(_{1}\) any free element of M. (Note that if x is a free element of M, the image of F under the map y → Φ(x, y) is the whole ring A; then proceed by induction to construct the bases (e \(_{n}\) ) and (f \(_{n}\) )). Deduce that if the basis (e \(_{n}\) ) of M is finite, E is the direct sum of M and N \(^{0}\) , and that one has M \(^{00}\) = M.
\item
  We keep the hypotheses of \(b\) , and suppose in addition that A satisfies condition \((\mathrm{N}_{s})\) and that \(d_{\Phi}\) is injective. Show then that there exists a basis \((e_{n})\) in E and a basis \((f_{n})\) in F such that \(\Phi(e_{i}, f_{j}) = \delta_{ij}\) . (Use \(b\) ) by determining alternately by induction \(e_{n}\) and \(f_{n}\) ).
\end{enumerate}

¶ 12) Let E be a real Hilbert space of countable type, \(\Phi(x, y)\) the scalar product in E. Show that there does not exist a system of two algebraic bases \((e_{\lambda})\) , \((f_{\mu})\) of the vector space E over R such that we have \(\Phi(e_{\lambda}, f_{\mu}) = \delta_{\lambda y}\) for every pair of indices. (First note that the index set of these bases would have the cardinality of the continuum (Topological vector spaces, chap. II, § 3, exerc. 15); then consider an orthonormal (countable) basis \((a_n)\) of E and note that the subspace spanned by the \(a_n\) is contained in the subspace generated by a countable subfamily of \((e_{\lambda})\) .
